\section{Онол, арга зүйн үндэслэл}

\subsection{Багцын өгөөж ба эрсдэл}


Энэхүү судалгаанд $N$ ширхэг хөрөнгөөс бүрдэх хөрөнгө оруулалтын багцыг авч үзье. $x_i$ нь $i$-р хөрөнгийн багцад эзлэх жин, $G_i$ нь тухайн хөрөнгийн өгөөж гэж тэмдэглэвэл багцын өгөөжийг
\[
r_p
=
\sum_{i=1}^{N}x_iG_i
\]
гэж тодорхойлно. Багцын хүлээгдэж буй өгөөж нь
\[
\mu_p
=
E[r_p]
=
E\left[
\sum_{i=1}^{N}x_iG_i
\right]
=
\sum_{i=1}^{N}x_iE[G_i]
\]
байна. Энд $E[G_i]$ нь $i$-р хөрөнгийн хүлээгдэж буй өгөөжийг илэрхийлнэ.

Багцын эрсдэлийг өгөөжийн вариацаар хэмжвэл
\[
\sigma_p^2
=
\operatorname{Var}(r_p)
=
\operatorname{Var}
\left[
\sum_{i=1}^{N}x_iG_i
\right]
=
\sum_{i=1}^{N}x_i^2\sigma_i^2
+
\sum_{i\neq j}x_ix_j\sigma_{ij},
\]
байна. Энд $\sigma_i^2=\operatorname{Var}(G_i)$ нь $i$-р
хөрөнгийн өгөөжийн вариац, харин
\[
\sigma_{ij}
=
\operatorname{Cov}(G_i,G_j)
\]
нь $i$ болон $j$-р хөрөнгийн өгөөжийн ковариац юм.

Багцын жингийн векторыг
\[
\mathbf{x}
=
(x_1,x_2,\ldots,x_N)^{\top},
\]
хөрөнгүүдийн өгөөжийн ковариацын матрицыг
\[
\Sigma
=
(\sigma_{ij})
=
\begin{pmatrix}
	\sigma_{11} & \sigma_{12} & \cdots & \sigma_{1N}\\
	\sigma_{21} & \sigma_{22} & \cdots & \sigma_{2N}\\
	\vdots      & \vdots      & \ddots & \vdots\\
	\sigma_{N1} & \sigma_{N2} & \cdots & \sigma_{NN}
\end{pmatrix}
\]
гэж тэмдэглэвэл багцын вариацыг матриц хэлбэрээр
\[
\sigma_p^2
=
\mathbf{x}^{\top}\Sigma\mathbf{x}
\]
гэж бичиж болно.

Иймээс багцын эрсдэл нь хөрөнгө тус бүрийн вариацаас гадна хөрөнгүүдийн өгөөж хоорондын ковариацаас хамаарна. Иймд ковариацын матрицыг хэрхэн үнэлэх нь багц сонголтын үр дүнд чухал нөлөөтэй.

Ковариацын матрицын үнэлэгчдийн харьцуулалтын үндсэн шалгуураар
түүврээс гадуурх багцын хэлбэлзлийг ашиглана
\citep{ledoitwolf2022,DeNardKostovic2026}.
Энэ нь хүлээгдэж буй өгөөжийн векторыг үнэлэх шаардлагагүйгээр
глобал хамгийн бага вариацтай багцын бодит хэрэгжилтийн
эрсдэлийг үнэлэх боломж олгоно.




\subsection{Глобал хамгийн бага вариацтай багц (Global Minimum Variance Portfolio, GMVP)}

Энэхүү судалгаанд ковариацын матрицын өөр өөр үнэлэгчийг багцын оновчлолд ашиглахад үүсэх ялгааг харьцуулахын тулд глобал хамгийн бага вариацтай багц (Global Minimum Variance Portfolio, GMVP)-ийг ашиглана. GMVP-ийн зорилго нь нийт хөрөнгө оруулалтыг багцад бүрэн хуваарилах нөхцөлд багцын вариацыг хамгийн бага байлгах хөрөнгө оруулалтын жинг тодорхойлоход оршино. GMVP-ийн оновчтой жинг тодорхойлоход хөрөнгүүдийн хүлээгдэж буй өгөөжийн үнэлгээ шаардлагагүй бөгөөд ковариацын матрицад тулгуурладаг.

Энэхүү судалгаанд богино арилжааг зөвшөөрөхгүй, хөрөнгө бүрийн багцад эзлэх жин сөрөг биш байх long-only GMVP загварыг ашиглана. Энэхүү хязгаарлалт нь богино арилжаа болон сөрөг жингээс үүдэх хөшүүрэгтэй байрлалыг хязгаарлахаас гадна ковариацын матрицын үнэлгээний алдаанаас шалтгаалан багцын жин хэт туйлшрах эрсдэлийг бууруулах боломжтой. \citet{JagannathanMa2003}-ийн судалгаанд багцын жинг сөрөг бишээр хязгаарлах нь үнэлсэн оновчтой багцын эрсдэлийг бууруулж болохыг харуулсан. Түүнчлэн ийм хязгаарлалт нь ковариацын матрицын зарим элементийг агшаасантай төстэй нөлөө үзүүлж болохыг тайлбарласан.

Гэхдээ long-only хязгаарлалт нь өөрөө ковариацын матрицын үнэлгээний алдааны нөлөөг бууруулж болох тул агшилтын үнэлэгчдийн хоорондын гүйцэтгэлийн ялгааг багасгаж болзошгүй. Иймээс энэхүү судалгаанд ковариацын матрицын үнэлэгчдийн үр дүнг ижил long-only оновчлолын нөхцөлийн хүрээнд харьцуулж, гарсан дүгнэлтийг уг хязгаарлалтын хүрээнд тайлбарлана.

$N$ ширхэг хөрөнгөөс бүрдэх багцын жингийн векторыг
\[
\boldsymbol{x}
=
(x_1,x_2,\ldots,x_N)^\top
\]
гэж тэмдэглэе. Мөн хөрөнгүүдийн өгөөжийн ковариацын матрицын үнэлэгчийг $\widehat{\Sigma}$ гэвэл багцын үнэлэгдсэн вариац
\[
\widehat{\sigma}_p^2
=
\boldsymbol{x}^{\top}
\widehat{\Sigma}
\boldsymbol{x}
\]
байна.

Тэгвэл long-only GMVP-ийн оновчлолын бодлогыг
\[
\begin{aligned}
	\min_{\boldsymbol{x}}\quad&
	\boldsymbol{x}^{\top}\widehat{\Sigma}\boldsymbol{x}\\
	\text{subject to}\quad&
	\boldsymbol{1}^{\top}\boldsymbol{x}=1,\\
	&x_i\geq 0,\qquad i=1,2,\ldots,N
\end{aligned}
\]
гэж тодорхойлно. Энд $\boldsymbol{1}$ нь бүх элемент нь нэгтэй $N$ хэмжээст вектор юм.

$\boldsymbol{1}^{\top}\boldsymbol{x}=1$ нөхцөл нь нийт хөрөнгө оруулалтыг багцад бүрэн хуваарилахыг, харин $x_i\geq 0$ нөхцөл нь хөрөнгө бүрийн жин сөрөг биш байхыг илэрхийлнэ. Эдгээр нөхцөлөөс
\[
0\leq x_i\leq 1,\qquad i=1,2,\ldots,N
\]
байх нь шууд мөрдөнө.

Long-only GMVP-ийн оновчтой жинг хязгаарлалттай квадратлаг оновчлолын бодлогын шийдээр тодорхойлно. Хэрэв $\widehat{\Sigma}$ нь эерэг хагас тодорхой матриц бол зорилгын функц гүдгэр байна. Мөн боломжит жингийн олонлог хоосон биш, хаалттай, хязгаарлагдмал тул оновчтой шийд оршино. Харин оновчтой шийд цорын ганц байх нь ковариацын матрицын шинж чанараас хамаарна.

Энэхүү судалгаанд түүврийн ковариацын матрицыг суурь үнэлэгч болгон авч, тогтмол вариацын, диагональ, тогтмол корреляцийн болон нэг хүчин зүйлийн зорилтот матрицад суурилсан шугаман агшилтын үнэлэгчид, мөн аналитик шугаман бус агшилтын үнэлэгчийг харьцуулна.

Эдгээр зургаан ковариацын матрицын үнэлэгчийг ерөнхийд нь
$\widehat{\Sigma}^{(k)}$, $k=1,2,\ldots,6$ гэж тэмдэглэе. Тэгвэл $k$-р үнэлэгчид харгалзах long-only GMVP-ийн оновчтой жинг
\[
\widehat{\boldsymbol{x}}^{(k)}
\in
\underset{
	\substack{
		\boldsymbol{x}\in\mathbb{R}^{N}\\
		\boldsymbol{1}^{\top}\boldsymbol{x}=1\\
		\boldsymbol{x}\geq\boldsymbol{0}
	}
}{\operatorname{arg\,min}}
\boldsymbol{x}^{\top}
\widehat{\Sigma}^{(k)}
\boldsymbol{x}
\]
гэж тодорхойлно.

Энд бүх үнэлэгчийн хувьд багцын оновчлолын зорилгын функцийн хэлбэр болон хөрөнгө оруулалтын жинд тавих хязгаарлалтууд ижил байх бөгөөд зөвхөн ковариацын матрицын үнэлэгч $\widehat{\Sigma}^{(k)}$ өөрчлөгдөнө.

Иймээс ковариацын матрицын үнэлэгчийн сонголтоос хамааран GMVP-ийн оновчтой жин, улмаар багцын эрсдэл болон бодит өгөөжийн үзүүлэлтүүд ялгаатай байж болно. Энэхүү судалгаанд эдгээр үнэлэгчээр байгуулсан long-only GMVP-үүдийн гүйцэтгэлийг ижил нөхцөлд харьцуулан үнэлнэ.





\subsection{Түүврийн ковариацын матриц}

Бодит нөхцөлд хөрөнгүүдийн өгөөжийн эх олонлогийн ковариацын матриц $\Sigma$ үл мэдэгдэх тул түүнийг ажиглагдсан өгөгдлөөс үнэлэх шаардлагатай. Энэхүү судалгаанд $n$ ажиглалт, $N$ хөрөнгөөс бүрдэх өгөөжийн ажиглалтын матрицыг
\[
Y
=
\begin{pmatrix}
	y_{11} & y_{12} & \cdots & y_{1N}\\
	y_{21} & y_{22} & \cdots & y_{2N}\\
	\vdots & \vdots & \ddots & \vdots\\
	y_{n1} & y_{n2} & \cdots & y_{nN}
\end{pmatrix}
\in\mathbb{R}^{n\times N}
\]
гэж тэмдэглэнэ. Энд $y_{ti}$ нь $t$-р хугацаанд ажиглагдсан $i$-р хөрөнгийн өгөөж, $n$ нь ажиглалтын тоо, $N$ нь хөрөнгийн тоо юм.

$i$-р хөрөнгийн дундаж өгөөжийг 
\[
\overline{Y}_i
=
\frac{1}{n}
\sum_{t=1}^{n}y_{ti},
\qquad
i=1,2,\ldots,N,
\]
гэж тооцно.

Тэгвэл $i$ болон $j$-р хөрөнгийн өгөөжийн түүврийн ковариацыг
\[
s_{ij}
=
\frac{1}{n}
\sum_{t=1}^{n}
(y_{ti}-\overline{Y}_i)
(y_{tj}-\overline{Y}_j)
\]
гэж тодорхойлно. Энд $1/n$ хуваарь ашигласан тул уг үнэлэгч нь
ерөнхий тохиолдолд хазайлтгүй түүврийн ковариацын үнэлэгч
($1/(n-1)$ хуваарьтай) биш болохыг анхаарна.
Эдгээр түүврийн ковариацуудаас бүрдэх
түүврийн ковариацын матриц нь
\[
S
=
(s_{ij})
=
\begin{pmatrix}
	s_{11} & s_{12} & \cdots & s_{1N}\\
	s_{21} & s_{22} & \cdots & s_{2N}\\
	\vdots & \vdots & \ddots & \vdots\\
	s_{N1} & s_{N2} & \cdots & s_{NN}
\end{pmatrix}
\]
байна. Энд $S$ нь $N\times N$ хэмжээст түүврийн ковариацын матриц бөгөөд эх олонлогийн ковариацын матриц $\Sigma$-г үнэлэхэд ашиглагдана.


\subsection{Шугаман агшилтын үнэлгээ (Linear Shrinkage Estimation)}

Шугаман агшилт (Linear Shrinkage) нь агшилтын үнэлгээний
хамгийн өргөн хэрэглэгддэг хэлбэрүүдийн нэг юм. Энэхүү аргын
үндсэн санаа нь түүврийн ковариацын матриц болон тодорхой
бүтэц бүхий зорилтот матрицыг гүдгэр хослол хэлбэрээр нэгтгэн
ковариацын матрицын үнэлэгчийг байгуулахад оршино. Шугаман
агшилтын үнэлэгчийг ерөнхий хэлбэрээр
\[
\widehat{\Sigma}_{LS}
=
(1-\alpha)S+\alpha T,
\qquad
0\leq\alpha\leq1,
\]
гэж тодорхойлно. Энд $S$ нь түүврийн ковариацын матриц,
$T$ нь зорилтот матриц, $\alpha$ нь агшилтын эрчим юм.
$\alpha=0$ үед $\widehat{\Sigma}_{LS}=S$, харин $\alpha=1$
үед $\widehat{\Sigma}_{LS}=T$ болно. Иймээс $\alpha$-ийн
утга ихсэх тусам зорилтот матриц $T$-д өгөх жин нэмэгдэж,
түүврийн ковариацын матриц $S$-д өгөх жин буурна.

Хэрэв зорилтот матриц $T$ нь эерэг тодорхой бөгөөд
$0<\alpha\leq1$ бол шугаман агшилтын үнэлэгч
$\widehat{\Sigma}_{LS}$ нь эерэг тодорхой байна. Энэ шинж нь
ковариацын матрицын тооцооллын тогтвортой байдлыг сайжруулахад
ач холбогдолтой.

Шугаман агшилтын арга нь түүврийн ковариацын матрицын бүх
элементэд нэг ижил агшилтын эрчим $\alpha$-г хэрэглэдэг.
Иймээс ковариацын матрицын бодит бүтэц нарийн төвөгтэй,
эсвэл хөрөнгүүдийн хоорондын хамаарал харилцан адилгүй үед
нэг ижил агшилтын эрчим хэрэглэх нь уг аргын уян хатан
байдлыг хязгаарлаж болно.

Шугаман агшилтын үнэлгээний үр дүнд зорилтот матрицын сонголт
чухал үүрэгтэй. Энэхүү судалгаанд тогтмол вариацын, диагональ,
тогтмол корреляцийн болон нэг хүчин зүйлийн зорилтот матрицыг
авч үзнэ. Эдгээр зорилтот матрицууд нь шугаман агшилтын
ерөнхий хэлбэр дэх $T$-ийн үүргийг гүйцэтгэх бөгөөд тэдгээрийн
ялгаа нь ковариацын матрицын бүтцэд тавих таамаглалд оршино.


% --------------------------------------------------
\subsubsection{Тогтмол вариацын зорилтот матриц}
% --------------------------------------------------

Тогтмол вариацын зорилтот матриц нь бүх хөрөнгийн
вариацыг нэг нийтлэг утгаар төлөөлүүлж, хөрөнгүүдийн
хоорондын ковариацыг тэг гэж үздэг.
Уг зорилтот матрицыг
\[
T_{\mathrm{var}}=\widehat{\mu}I
\]
гэж тодорхойлно. Энд $I$ нь $N\times N$
хэмжээст нэгж матриц, харин
\[
\widehat{\mu}
=
\frac{1}{N}\sum_{i=1}^{N}s_{ii}
=
\frac{\operatorname{tr}(S)}{N}
\]
нь хөрөнгүүдийн өгөөжийн түүврийн
вариацуудын дундаж юм. Иймээс
\[
T_{\mathrm{var}}
=
\begin{pmatrix}
	\widehat{\mu}&0&\cdots&0\\
	0&\widehat{\mu}&\cdots&0\\
	\vdots&\vdots&\ddots&\vdots\\
	0&0&\cdots&\widehat{\mu}
\end{pmatrix}
\]
хэлбэртэй байна \citep{ledoitwolf2004b}.

Тогтмол вариацын зорилтот матрицад
суурилсан шугаман агшилтын үнэлэгчийг
\[
\widehat{\Sigma}_{\mathrm{var}}
=
(1-\widehat{\alpha}_{\mathrm{var}})S
+
\widehat{\alpha}_{\mathrm{var}}T_{\mathrm{var}},
\qquad
0\leq\widehat{\alpha}_{\mathrm{var}}\leq1
\]
гэж тодорхойлно.

Агшилтын эрчмийн онолын үндэслэлийг
Ledoit болон Wolf (2004)-ийн үр дүнд
тулгуурлан дараах хэмжигдэхүүнүүдийг авч үзье:
\[
\begin{aligned}
	a&=E\left[\|\Sigma-T\|^2\right]=\|\Sigma-T\|^2,\\
	b&=E\left[\|S-\Sigma\|^2\right],\\
	d&=E\left[\|S-T\|^2\right].
\end{aligned}
\]

\medskip
\noindent
\textbf{Лемма 1 (Ledoit--Wolf, 2004).}
Хэрэв $E[S]=\Sigma$ бөгөөд квадратын математик дунджууд төгсгөлөг бол
\[
a+b=d
\]
тэнцэтгэл биелнэ.

\medskip
\noindent
\textit{Баталгаа.}
Математик дундаж ба нормын чанараас
\[
\begin{aligned}
	d
	&=E\left[\|S-T\|^2\right]=E\left[\|S-\Sigma+\Sigma-T\|^2\right]\\
	&=E\left[\|S-\Sigma\|^2\right]
	+\|\Sigma-T\|^2+
	2E\left[
	\langle S-\Sigma,\Sigma-T\rangle
	\right].
\end{aligned}
\]

$E[S]=\Sigma$ бөгөөд $T$ нь
санамсаргүй бус матриц тул
\[
E\left[
\langle S-\Sigma,\Sigma-T\rangle
\right]
=
\langle E[S-\Sigma],\Sigma-T\rangle
=0.
\]

Иймээс
\[
d=b+a.
\]
Лемма батлагдав. \hfill $\square$

\medskip
\noindent
\textbf{Теорем 1 (Ledoit--Wolf, 2004).}
Дараах оновчлолын бодлогыг авч үзье:
\[
\min_{0\leq\alpha\leq1}
E\left[
\|\Sigma^*(\alpha)-\Sigma\|^2
\right],
\]
энд
\[
\Sigma^*(\alpha)
=
(1-\alpha)S+\alpha T.
\]

Хэрэв $d>0$ бол оновчтой
агшилтын эрчим
\[
\alpha^*=\frac{b}{d}
\]
бөгөөд харгалзах үнэлэгч
\[
\Sigma^*
=
\frac{a}{d}S
+
\frac{b}{d}T
\]
байна.

\medskip
\noindent
\textit{Баталгаа.}
Агшилтын үнэлэгчийн алдааг
\[
\Sigma^*(\alpha)-\Sigma
=
(1-\alpha)(S-\Sigma)
+\alpha(T-\Sigma)
\]
гэж бичиж болно.

Одоо нормын квадратын томьёо болон математик дунджийн шугаман чанарыг хэрэглэнэ.
\[
\begin{aligned}
	\|\Sigma^*(\alpha)-\Sigma\|^2&=\|(1-\alpha)(S-\Sigma)
	+\alpha(T-\Sigma)\|^2\\
	&=(1-\alpha)^2\|S-\Sigma\|^2+\alpha^2\|T-\Sigma\|^2+2\alpha(1-\alpha)
	\langle S-\Sigma,T-\Sigma\rangle.
\end{aligned}
\]

Хоёр талаас математик дундаж авбал:
\[
\begin{aligned}
	E\left[\|\Sigma^*(\alpha)-\Sigma\|^2\right]&=(1-\alpha)^2E\left[\|S-\Sigma\|^2\right]+\alpha^2E\left[\|T-\Sigma\|^2\right]+2\alpha(1-\alpha)E\left[\langle S-\Sigma,T-\Sigma\rangle\right]\\
	&=
	(1-\alpha)^2b+\alpha^2a.
\end{aligned}
\]

Иймээс зорилгын функц
\[
R(\alpha)
=
(1-\alpha)^2b+\alpha^2a
\]
болно.

Алдааг хамгийн бага болгох \(\alpha\)-г олох нь бидний зорилго учраас уламжлал аваад тэгтэй тэнцүүлнэ.
\[
R'(\alpha)
=
-2(1-\alpha)b+2\alpha a=0.
\]

Лемма 1-ээс $a+b=d$ тул
\[
\alpha^*=\frac{b}{d}, \qquad 
1-\alpha^*=\frac{a}{d}.
\]

Иймээс
\[
\Sigma^*
=
\frac{a}{d}S
+
\frac{b}{d}T.
\]

$R''(\alpha)=2d>0$ учраас
энэ нь минимум шийд юм.

Теорем батлагдав. \hfill $\square$

\medskip



Теорем 1-ээс оновчтой агшилтын
эрчим нь түүврийн ковариацын матрицын
үнэлгээний алдаа болон зорилтот
матрицаас хазайх хэмжээнээс
хамаарах нь харагдана.

Гэвч $b$ болон $d$ нь үл мэдэгдэх
бодит ковариацын матрицаас хамаардаг
тул тэдгээрийг өгөгдлөөс үнэлэх
шаардлагатай. Ledoit болон Wolf (2004)-ийн
оновчтой агшилтын
эрчмийг өгөгдлөөс үнэлэх аргыг
боловсруулсан.

Түүврийн ковариацын матриц болон
зорилтот матрицын хоорондын
квадрат Фробениусын зайг
\[
\widehat{d}_{\mathrm{var}}
=
\|S-T_{\mathrm{var}}\|_F^2
=
\frac{1}{N}\sum_{i=1}^{N}\sum_{j=1}^{N}
(s_{ij}-t_{ij})^2
\]
гэж тодорхойлно.

Түүврийн ковариацын матрицын хүлээгдэж буй
квадрат алдааг нормчилсон Фробениусын нормын
тодорхойлолтоос
\[
b_{var}=E\left[\|S-\Sigma\|^2\right]=\frac{1}{N}
\sum_{i=1}^{N}\sum_{j=1}^{N}
E\left[(s_{ij}-\sigma_{ij})^2\right].
\]

Хэрэв $E[S]=\Sigma$ бол $E\left[(s_{ij}-\sigma_{ij})^2\right]
=
\operatorname{Var}(s_{ij}).$

Иймээс
\[
b_{var}
=
\frac{1}{N}
\sum_{i=1}^{N}\sum_{j=1}^{N}
\operatorname{Var}(s_{ij}).
\]

Одоо
\[
s_{ij}=\frac{1}{n}\sum_{t=1}^{n}\{(y_{ti}-\overline{Y}_i)
(y_{tj}-\overline{Y}_j)\}
\]
гэж тэмдэглэе.

Ажиглалтууд хугацааны хувьд
харилцан хамааралгүй, ижил тархалттай
бол
\[
\operatorname{Var}(s_{ij})\approx
\frac{1}{n^2}
\sum_{t=1}^{n}\operatorname{Var}\left((y_{ti}-\overline{Y}_i)
(y_{tj}-\overline{Y}_j)\right)=
\frac{1}{n}\operatorname{Var}\left((y_{1i}-\overline{Y}_i)
(y_{1j}-\overline{Y}_j)\right).
\]

Тэгвэл
\[
b_{var}=\frac{1}{N}
\sum_{i=1}^{N}\sum_{j=1}^{N}
\frac{1}{n}\operatorname{Var}\left((y_{1i}-\overline{Y}_i)
(y_{1j}-\overline{Y}_j)\right)=
\frac{1}{nN}
\sum_{i=1}^{N}\sum_{j=1}^{N}
\operatorname{Var}\left((y_{1i}-\overline{Y}_i)
(y_{1j}-\overline{Y}_j)\right).
\]

Энд $\operatorname{Var}\left((y_{1i}-\overline{Y}_i)
(y_{1j}-\overline{Y}_j)\right)$ нь
үл мэдэгдэх тул түүний үнэлгээг $\frac{1}{n}
\sum_{t=1}^{n}((y_{ti}-\overline{Y}_i)
(y_{tj}-\overline{Y}_j)-s_{ij})^2$ байгуулвал
\[
\widehat{b}_{var}=
\frac{1}{nN}
\sum_{i=1}^{N}\sum_{j=1}^{N}
\left[
\frac{1}{n}\sum_{t=1}^{n}
\left\{(y_{ti}-\overline{Y}_i)
(y_{tj}-\overline{Y}_j)-s_{ij}\right\}^2
\right]=
\]
\[=
\frac{1}{nN}
\left[
\frac{1}{n}
\sum_{t=1}^{n}
\sum_{i=1}^{N}\sum_{j=1}^{N}
\left\{(y_{ti}-\overline{Y}_i)
(y_{tj}-\overline{Y}_j)-s_{ij}\right\}^2
\right].
\]
Иймээс агшилтын эрчмийн
практик үнэлгээг
\[
\widehat{\alpha}_{\mathrm{var}}
=
\max\left\{
0,\,
\min\left\{
1,\,
\frac{\widehat{b}_{var}}
{\widehat{d}_{\mathrm{var}}}
\right\}
\right\}
\]
гэж байгуулна \citep{ledoitwolf2004b}.


\paragraph{Тооцоолох алгоритм.}

Тогтмол вариацын зорилтот матрицад суурилсан
агшилтын ковариацын матрицыг дараах
дарааллаар тооцоолно.

\begin{enumerate}
	\item Түүврийн ковариацын матрицыг байгуулах.
	
	$n$ ажиглалт, $N$ хөрөнгийн өгөөжийн
	өгөгдлөөс
	\[
	s_{ij}
	=
	\frac{1}{n}
	\sum_{t=1}^{n}
	(y_{ti}-\overline{Y}_i)
	(y_{tj}-\overline{Y}_j)
	\]
	гэж тооцож, $S=(s_{ij})$ матрицыг байгуулна.
	
	\item Тогтмол вариацын зорилтот
	матрицыг байгуулах.
	
	\[
	\widehat{\mu}
	=
	\frac{1}{N}\sum_{i=1}^{N}s_{ii},
	\qquad
	T_{\mathrm{var}}=\widehat{\mu}I.
	\]
	
	\item Ковариацын матрицын
	квадрат алдааг үнэлэх.
	\[
	\widehat{b}_{var}
	=
	\frac{1}{nN}
	\sum_{i=1}^{N}\sum_{j=1}^{N}
	\left[
	\frac{1}{n}
	\sum_{t=1}^{n}
	\{(y_{ti}-\overline{Y}_i)
	(y_{tj}-\overline{Y}_j)-s_{ij}\}^2
	\right].
	\]
	
	\item Зорилтот матрицаас хазайх
	квадрат зайг тооцох.
	
	\[
	\widehat{d}_{\mathrm{var}}
	=
	\frac{1}{N}
	\sum_{i=1}^{N}\sum_{j=1}^{N}
	(s_{ij}-t_{ij})^2.
	\]
	
	\item Агшилтын эрчмийг үнэлэх.
	
	$\widehat{d}_{\mathrm{var}}>0$
	үед
	\[
	\widehat{\alpha}_{\mathrm{var}}
	=
	\max\left\{
	0,\,
	\min\left\{
	1,\,
	\frac{\widehat{b}_{var}}
	{\widehat{d}_{\mathrm{var}}}
	\right\}
	\right\}.
	\]
	
	\item Агшилтын ковариацын
	матрицыг байгуулах.
	
	\[
	\boxed{
		\widehat{\Sigma}_{\mathrm{var}}
		=
		(1-\widehat{\alpha}_{\mathrm{var}})S
		+
		\widehat{\alpha}_{\mathrm{var}}
		T_{\mathrm{var}}
	}.
	\]
\end{enumerate}

% --------------------------------------------------
\subsubsection{Диагональ зорилтот матриц}
% --------------------------------------------------

Диагональ зорилтот матриц нь хөрөнгө тус бүрийн өгөөжийн
түүврийн вариацыг хадгалж, харин хөрөнгүүдийн өгөөж хоорондын
ковариацыг тэг гэж үзнэ. Уг зорилтот матрицыг

\[
T_{\mathrm{diag}}
=
\operatorname{diag}(S)
=
\operatorname{diag}
(s_{11},s_{22},\ldots,s_{NN})
\]

гэж тодорхойлно. Матриц хэлбэрээр

\[
T_{\mathrm{diag}}
=
\begin{pmatrix}
	s_{11} & 0      & \cdots & 0\\
	0      & s_{22} & \cdots & 0\\
	\vdots & \vdots & \ddots & \vdots\\
	0      & 0      & \cdots & s_{NN}
\end{pmatrix}.
\]

Өөрөөр хэлбэл,

\[
(T_{\mathrm{diag}})_{ij}
=
\begin{cases}
	s_{ii}, & i=j,\\
	0,      & i\neq j.
\end{cases}
\]

Иймээс диагональ зорилтот матриц нь хөрөнгө тус бүрийн
өгөөжийн түүврийн вариацыг өөрчлөхгүй хадгалах боловч
хөрөнгүүдийн өгөөж хоорондын ковариацыг тэг болгон
хялбаршуулдаг.

Диагональ зорилтот матрицад суурилсан шугаман агшилтын
үнэлэгчийг

\[
\widehat{\Sigma}_{\mathrm{diag}}
=
(1-\widehat{\alpha}_{\mathrm{diag}})S
+
\widehat{\alpha}_{\mathrm{diag}}T_{\mathrm{diag}},
\qquad
0\leq\widehat{\alpha}_{\mathrm{diag}}\leq1,
\]

гэж тодорхойлно. Энд
$\widehat{\alpha}_{\mathrm{diag}}$ нь диагональ зорилтот
матрицад харгалзах агшилтын эрчим юм.



\paragraph{Агшилтын эрчмийн үнэлгээ.}

Диагональ зорилтот матриц нь түүврийн ковариацын
матрицын диагональ элементүүдийг өөрчлөхгүй
хадгалдаг тул агшилт нь зөвхөн диагоналийн бус
элементүүдэд үйлчилнэ.

Иймээс агшилтын эрчмийг үнэлэхдээ
диагоналийн бус элементүүдийн квадрат
алдааг авч үзнэ.

Түүврийн ковариацын матриц болон
диагональ зорилтот матрицын хоорондын
нормчилсон квадрат Фробениусын зайг
\[
\widehat d_{\mathrm{diag}}
=
\|S-T_{\mathrm{diag}}\|^2
=
\frac{1}{N}
\sum_{\substack{i,j=1\\i\neq j}}^{N}
s_{ij}^{2}
\]
гэж тодорхойлно.

Учир нь $S$ болон $T_{\mathrm{diag}}$
матрицын диагональ элементүүд ижил тул
тэдгээрийн ялгавар диагональ дээр тэг байна.

Өмнөх тогтмол вариацын зорилтот матрицын хэсэгт үзсэнтэй адилтаар $\widehat{b}_{diag}$-г дараах байдлаар тодорхойлно. 
\[
\widehat b_{\mathrm{diag}}
=
\frac{1}{nN}
\sum_{\substack{i,j=1\\i\neq j}}^{N}
\left[
\frac{1}{n}
\sum_{t=1}^{n}
(z_{tij}-s_{ij})^2
\right].
\]
Иймээс диагональ зорилтот матрицад
харгалзах агшилтын эрчмийн
практик үнэлгээг
\[
\widehat{\alpha}_{\mathrm{diag}}
=
\max\left\{
0,\,
\min\left\{
1,\,
\frac{\widehat b_{\mathrm{diag}}}
{\widehat d_{\mathrm{diag}}}
\right\}
\right\}
\]
гэж байгуулна.

Эцэст нь агшилтын ковариацын
матрицыг
\[
\widehat{\Sigma}_{\mathrm{diag}}
=
(1-\widehat{\alpha}_{\mathrm{diag}})S
+
\widehat{\alpha}_{\mathrm{diag}}
T_{\mathrm{diag}}
\]
гэж тооцно.

\paragraph{Тооцоолох алгоритм.}

Диагональ зорилтот матрицад суурилсан
агшилтын ковариацын матрицыг дараах
дарааллаар тооцоолно.

\begin{enumerate}
	\item Диагональ зорилтот матрицыг байгуулах.
	
	Өмнөх хэсэгт тооцсон түүврийн ковариацын
	матриц $S=(s_{ij})$-ийг ашиглан
	\[
	T_{\mathrm{diag}}
	=
	\operatorname{diag}
	(s_{11},s_{22},\ldots,s_{NN})
	\]
	гэж байгуулна.
	\item Диагоналийн бус элементүүдийн
	квадрат алдааг үнэлэх.
	\[
	\widehat b_{\mathrm{diag}}
	=
	\frac{1}{nN}
	\sum_{\substack{i,j=1\\i\neq j}}^{N}
	\left[
	\frac{1}{n}
	\sum_{t=1}^{n}
	(z_{tij}-s_{ij})^2
	\right].
	\]
	\item Зорилтот матрицаас хазайх
	квадрат зайг тооцох.
	
	\[
	\widehat d_{\mathrm{diag}}
	=
	\frac{1}{N}
	\sum_{\substack{i,j=1\\i\neq j}}^{N}
	s_{ij}^{2}.
	\]
	
	\item Агшилтын эрчмийг үнэлэх.
	
	$\widehat d_{\mathrm{diag}}>0$ үед
	\[
	\widehat{\alpha}_{\mathrm{diag}}
	=
	\max\left\{
	0,\,
	\min\left\{
	1,\,
	\frac{\widehat b_{\mathrm{diag}}}
	{\widehat d_{\mathrm{diag}}}
	\right\}
	\right\}.
	\]
	
	\item Агшилтын ковариацын
	матрицыг байгуулах.
	
	\[
	\widehat{\Sigma}_{\mathrm{diag}}
	=
	(1-\widehat{\alpha}_{\mathrm{diag}})S
	+
	\widehat{\alpha}_{\mathrm{diag}}
	T_{\mathrm{diag}}.
	\]
\end{enumerate}


% --------------------------------------------------
\subsubsection{Тогтмол корреляцийн зорилтот матриц}
% --------------------------------------------------

Тогтмол корреляцийн зорилтот матриц нь хөрөнгө тус бүрийн
өгөөжийн түүврийн вариацыг хадгалж, харин хөрөнгүүдийн
хоорондын корреляцыг тэдгээрийн дундаж утгаар
төлөөлүүлэн байгуулдаг зорилтот матриц юм
\citep{ledoitwolf2004a}.

$i$ болон $j$ дугаар хөрөнгийн өгөөжийн
түүврийн корреляцийн коэффициентийг
\[
r_{ij}
=
\frac{s_{ij}}{\sqrt{s_{ii}s_{jj}}}
\]
гэж тодорхойлно \citep{ledoitwolf2004a}.

Хөрөнгүүдийн хоорондын түүврийн
корреляцуудын дундажийг
\[
\bar r
=
\frac{2}{N(N-1)}
\sum_{i=1}^{N-1}
\sum_{j=i+1}^{N}r_{ij}
\]
гэж тооцно.

Тогтмол корреляцийн зорилтот матрицын
элементүүдийг
\[
(T_{\mathrm{corr}})_{ij}
=
t_{ij}
=
\begin{cases}
	s_{ii}, & i=j,\\
	\bar r\sqrt{s_{ii}s_{jj}}, & i\neq j
\end{cases}
\]
гэж тодорхойлно.

Иймээс уг зорилтот матриц нь хөрөнгө тус бүрийн
өгөөжийн түүврийн вариацыг өөрчлөхгүй хадгалж,
харин хөрөнгүүдийн хоорондын корреляцыг
нэг ижил дундаж утгаар орлуулна.

Тогтмол корреляцийн зорилтот матрицад
суурилсан шугаман агшилтын үнэлэгчийг
\[
\widehat{\Sigma}_{\mathrm{corr}}
=
(1-\widehat{\alpha}_{\mathrm{corr}})S
+
\widehat{\alpha}_{\mathrm{corr}}T_{\mathrm{corr}},
\qquad
0\leq\widehat{\alpha}_{\mathrm{corr}}\leq1
\]
гэж тодорхойлно. Энд
$\widehat{\alpha}_{\mathrm{corr}}$ нь тогтмол
корреляцийн зорилтот матрицад харгалзах
агшилтын эрчим юм.

\paragraph{Агшилтын эрчмийн үнэлгээ.}

Тогтмол корреляцийн зорилтот матриц нь
түүврийн ковариацын матрицын элементүүдээс
хамааран байгуулагддаг тул агшилтын
эрчмийг үнэлэхдээ түүврийн ковариацын
матриц болон зорилтот матрицын
хам хэлбэлзлийг харгалзан үзнэ.

Түүврийн ковариацын матриц болон
тогтмол корреляцийн зорилтот матрицын
хоорондын нормчилсон квадрат
Фробениусын зайг
\[
\widehat d_{\mathrm{corr}}
=
\|S-T_{\mathrm{corr}}\|^2
=
\frac{1}{N}
\sum_{i=1}^{N}\sum_{j=1}^{N}
(s_{ij}-t_{ij})^2
=
\frac{1}{N}
\sum_{\substack{i,j=1\\i\neq j}}^{N}
\left(
s_{ij}-\bar r\sqrt{s_{ii}s_{jj}}
\right)^2
\]
гэж тодорхойлно.

Түүврийн ковариацын матрицын
хүлээгдэж буй квадрат алдааг өмнөх хэсгүүдэд үзсэнтэй адилаар 
\[
\widehat b_{\mathrm{corr}}
=
\frac{1}{nN}
\sum_{i=1}^{N}\sum_{j=1}^{N}
\left[
\frac{1}{n}
\sum_{t=1}^{n}
\left\{
(y_{ti}-\overline{Y}_i)
(y_{tj}-\overline{Y}_j)-s_{ij}
\right\}^2
\right]
\]
гэж бичигдэнэ.

Тогтмол корреляцийн зорилтот матриц
нь түүврийн ковариацын матрицаас
хамаардаг учраас тэдгээрийн
харгалзах элементүүдийн
асимптотик ковариацыг тооцох
шаардлагатай.

Энэ хэмжигдэхүүний нормчилсон
үнэлгээг $\widehat c_{\mathrm{corr}}$
гэж тэмдэглэе.

Эхлээд
\[
\widehat{\vartheta}_{ii,ij}
=
\frac{1}{n}
\sum_{t=1}^{n}
\left\{
(y_{ti}-\overline{Y}_i)^2-s_{ii}
\right\}
\left\{
(y_{ti}-\overline{Y}_i)
(y_{tj}-\overline{Y}_j)-s_{ij}
\right\}
\]
болон
\[
\widehat{\vartheta}_{jj,ij}
=
\frac{1}{n}
\sum_{t=1}^{n}
\left\{
(y_{tj}-\overline{Y}_j)^2-s_{jj}
\right\}
\left\{
(y_{ti}-\overline{Y}_i)
(y_{tj}-\overline{Y}_j)-s_{ij}
\right\}
\]
гэж тодорхойлно.

Эдгээр нь түүврийн вариац болон
ковариацын элементүүдийн
асимптотик ковариацыг үнэлэхэд
ашиглагдана.

Ledoit болон Wolf (2004) нь түүврийн ковариацын
матриц болон тогтмол корреляцийн зорилтот матрицын
асимптотик хам хэлбэлзлийг тооцох аргыг
боловсруулсан \citep{ledoitwolf2004a}.
Тэдгээрийн Appendix B-д өгсөн үр дүнд тулгуурлан
\[
\widehat c_{\mathrm{corr}}
=
\frac{1}{nN}
\Bigg[
\sum_{i=1}^{N}
\widehat{\operatorname{Var}}
\left((y_{1i}-\overline{Y}_i)^2\right)+
\frac{\bar r}{2}
\sum_{\substack{i,j=1\\i\neq j}}^{N}
\left(
\sqrt{\frac{s_{jj}}{s_{ii}}}
\widehat{\vartheta}_{ii,ij}
+
\sqrt{\frac{s_{ii}}{s_{jj}}}
\widehat{\vartheta}_{jj,ij}
\right)
\Bigg]
\]
гэж үнэлнэ.

Энд
\[
\widehat{\operatorname{Var}}
\left((y_{1i}-\overline{Y}_i)^2\right)
=
\frac{1}{n}
\sum_{t=1}^{n}
\left\{
(y_{ti}-\overline{Y}_i)^2-s_{ii}
\right\}^2
\]
юм.

Иймээс тогтмол корреляцийн
зорилтот матрицад харгалзах
агшилтын эрчмийн практик үнэлгээг
\[
\widehat{\alpha}_{\mathrm{corr}}
=
\max\left\{
0,\,
\min\left\{
1,\,
\frac{
	\widehat b_{\mathrm{corr}}
	-
	\widehat c_{\mathrm{corr}}
}{
	\widehat d_{\mathrm{corr}}
}
\right\}
\right\}
\]
гэж байгуулна \citep{ledoitwolf2004a}.

Энд $\widehat b_{\mathrm{corr}}$ нь
түүврийн ковариацын матрицын
үнэлгээний квадрат алдааны
асимптотик ойролцоо үнэлгээ,
$\widehat c_{\mathrm{corr}}$ нь
түүврийн ковариацын матриц болон
зорилтот матрицын харгалзах
элементүүдийн асимптотик
ковариацуудын нийлбэрийн
нормчилсон үнэлгээ,
$\widehat d_{\mathrm{corr}}$ нь
хоёр матрицын хоорондын
нормчилсон квадрат Фробениусын
зай юм.

Эцэст нь агшилтын ковариацын
матрицыг
\[
\widehat{\Sigma}_{\mathrm{corr}}
=
(1-\widehat{\alpha}_{\mathrm{corr}})S
+
\widehat{\alpha}_{\mathrm{corr}}
T_{\mathrm{corr}}
\]
гэж тооцно.

\paragraph{Тооцоолох алгоритм.}

Тогтмол корреляцийн зорилтот матрицад
суурилсан агшилтын ковариацын матрицыг
дараах дарааллаар тооцоолно.

\begin{enumerate}
	\item Тогтмол корреляцийн зорилтот
	матрицыг байгуулах.
	
	Өмнөх хэсэгт тооцсон түүврийн
	ковариацын матриц $S=(s_{ij})$-ийг
	ашиглан
	\[
	r_{ij}
	=
	\frac{s_{ij}}{\sqrt{s_{ii}s_{jj}}},
	\qquad
	\bar r
	=
	\frac{2}{N(N-1)}
	\sum_{i=1}^{N-1}
	\sum_{j=i+1}^{N}r_{ij}
	\]
	гэж тооцно.
	
	Дараа нь
	\[
	t_{ij}
	=
	\begin{cases}
		s_{ii}, & i=j,\\
		\bar r\sqrt{s_{ii}s_{jj}}, & i\neq j
	\end{cases}
	\]
	гэж тодорхойлж,
	$T_{\mathrm{corr}}=(t_{ij})$
	матрицыг байгуулна.
	
	\item Ковариацын матрицын
	квадрат алдааг үнэлэх.
	
	\[
	\widehat b_{\mathrm{corr}}
	=
	\frac{1}{nN}
	\sum_{i=1}^{N}\sum_{j=1}^{N}
	\left[
	\frac{1}{n}
	\sum_{t=1}^{n}
	\left\{
	(y_{ti}-\overline{Y}_i)
	(y_{tj}-\overline{Y}_j)-s_{ij}
	\right\}^2
	\right]
	\]
	гэж тооцно.
	
	\item Түүврийн ковариацын матриц
	болон зорилтот матрицын
	хам хэлбэлзлийг үнэлэх.
	
	\[
	\widehat{\vartheta}_{ii,ij}
	=
	\frac{1}{n}
	\sum_{t=1}^{n}
	\left\{
	(y_{ti}-\overline{Y}_i)^2-s_{ii}
	\right\}
	\left\{
	(y_{ti}-\overline{Y}_i)
	(y_{tj}-\overline{Y}_j)-s_{ij}
	\right\}
	\]
	болон
	\[
	\widehat{\vartheta}_{jj,ij}
	=
	\frac{1}{n}
	\sum_{t=1}^{n}
	\left\{
	(y_{tj}-\overline{Y}_j)^2-s_{jj}
	\right\}
	\left\{
	(y_{ti}-\overline{Y}_i)
	(y_{tj}-\overline{Y}_j)-s_{ij}
	\right\}
	\]
	гэж тооцно.
	
	Дараа нь
	\[
	\widehat c_{\mathrm{corr}}
	=
	\frac{1}{nN}
	\Bigg[
	\sum_{i=1}^{N}
	\widehat{\operatorname{Var}}
	\left((y_{1i}-\overline{Y}_i)^2\right)+
	\frac{\bar r}{2}
	\sum_{\substack{i,j=1\\i\neq j}}^{N}
	\left(
	\sqrt{\frac{s_{jj}}{s_{ii}}}
	\widehat{\vartheta}_{ii,ij}
	+
	\sqrt{\frac{s_{ii}}{s_{jj}}}
	\widehat{\vartheta}_{jj,ij}
	\right)
	\Bigg]
	\]
	гэж үнэлнэ.
	
	\item Зорилтот матрицаас хазайх
	квадрат зайг тооцох.
	
	\[
	\widehat d_{\mathrm{corr}}
	=
	\frac{1}{N}
	\sum_{\substack{i,j=1\\i\neq j}}^{N}
	(s_{ij}-t_{ij})^2.
	\]
	
	\item Агшилтын эрчмийг үнэлэх.
	
	$\widehat d_{\mathrm{corr}}>0$
	үед
	\[
	\widehat{\alpha}_{\mathrm{corr}}
	=
	\max\left\{
	0,\,
	\min\left\{
	1,\,
	\frac{
		\widehat b_{\mathrm{corr}}
		-
		\widehat c_{\mathrm{corr}}
	}{
		\widehat d_{\mathrm{corr}}
	}
	\right\}
	\right\}.
	\]
	
	\item Агшилтын ковариацын
	матрицыг байгуулах.
	
	\[
	\widehat{\Sigma}_{\mathrm{corr}}
	=
	(1-\widehat{\alpha}_{\mathrm{corr}})S
	+
	\widehat{\alpha}_{\mathrm{corr}}
	T_{\mathrm{corr}}.
	\]
\end{enumerate}

% --------------------------------------------------
\subsubsection{Нэг хүчин зүйлийн зорилтот матриц}
% --------------------------------------------------

Нэг хүчин зүйлийн зорилтот матриц нь хөрөнгүүдийн өгөөжийн
хоорондын хамаарлыг зах зээлийн нийтлэг хүчин зүйлээр
тайлбарлах нэг хүчин зүйлийн загварт (Single-Index Model)
суурилдаг. Уг загварт хөрөнгийн өгөөжийг зах зээлийн
нийтлэг хүчин зүйлтэй холбоотой хэсэг болон тухайн
хөрөнгийн өөрийн онцлогоос хамаарах үлдэгдэл хэсэгт
хуваан авч үздэг \citep{ledoitwolf2003}.

Энэхүү судалгаанд зах зээлийн нийтлэг хүчин зүйлийг
S\&P 500 индексийн өгөөжөөр төлөөлүүлнэ.

$i$-р хөрөнгийн $t$-р хугацаан дахь өгөөжийг
\[
y_{ti}
=
a_i+\beta_i R_{m,t}+\varepsilon_{ti}
\]
гэж загварчилна. Энд $R_{m,t}$ нь $t$-р хугацаан дахь
зах зээлийн өгөөж, $a_i$ нь регрессийн тогтмол
коэффициент, $\beta_i$ нь $i$-р хөрөнгийн өгөөжийн
зах зээлийн өөрчлөлтөд мэдрэг байдлыг илэрхийлэх
коэффициент, $\varepsilon_{ti}$ нь тухайн хөрөнгийн
үлдэгдэл өгөөж юм.

Нэг хүчин зүйлийн загварын регрессийн коэффициентийг
\[
\widehat{\beta}_i
=
\frac{
	\sum_{t=1}^{n}
	(y_{ti}-\overline{Y}_i)
	(R_{m,t}-\overline{R}_m)
}{
	\sum_{t=1}^{n}
	(R_{m,t}-\overline{R}_m)^2
}
\]
гэж үнэлнэ. Энд
\[
\overline{R}_m
=
\frac{1}{n}
\sum_{t=1}^{n}R_{m,t}
\]
нь зах зээлийн өгөөжийн түүврийн дундаж юм.

Регрессийн тогтмол коэффициентийг
\[
\widehat{a}_i
=
\overline{Y}_i
-
\widehat{\beta}_i\overline{R}_m
\]
гэж тооцно.

Тэгвэл $i$-р хөрөнгийн $t$-р хугацаан дахь
регрессийн үлдэгдэл
\[
\widehat{\varepsilon}_{ti}
=
y_{ti}
-
\widehat{a}_i
-
\widehat{\beta}_iR_{m,t}
\]
болно. Үүнийг
\[
\widehat{\varepsilon}_{ti}
=
(y_{ti}-\overline{Y}_i)
-
\widehat{\beta}_i
(R_{m,t}-\overline{R}_m)
\]
гэж илэрхийлж болно.

Нэг хүчин зүйлийн загварын үндсэн таамаглалаар
үлдэгдэл өгөөжүүд нь зах зээлийн хүчин зүйлтэй
хамааралгүй, мөн өөр өөр хөрөнгүүдийн үлдэгдэл
өгөөжүүд хоорондоо хамааралгүй байна. Өөрөөр хэлбэл,
\[
\operatorname{Cov}(R_m,\varepsilon_i)=0,
\qquad i=1,2,\ldots,N,
\]
ба
\[
\operatorname{Cov}(\varepsilon_i,\varepsilon_j)=0,
\qquad i\neq j.
\]

Эдгээр таамаглалын дор $i$ болон $j$-р хөрөнгийн
өгөөжийн ковариацыг
\[
\operatorname{Cov}(Y_i,Y_j)
=
\operatorname{Cov}
(a_i+\beta_iR_m+\varepsilon_i,
a_j+\beta_jR_m+\varepsilon_j)
\]
гэж илэрхийлнэ.

Ковариацын чанарыг ашиглавал
\[
\operatorname{Cov}(Y_i,Y_j)
=
\beta_i\beta_j\operatorname{Var}(R_m),
\qquad i\neq j,
\]
болно.

Харин $i$-р хөрөнгийн өгөөжийн вариац
\[
\operatorname{Var}(Y_i)
=
\beta_i^2\operatorname{Var}(R_m)
+
\operatorname{Var}(\varepsilon_i)
\]
болно.

Зах зээлийн өгөөжийн түүврийн вариацыг
\[
\widehat{\sigma}_m^2
=
\frac{1}{n}
\sum_{t=1}^{n}
(R_{m,t}-\overline{R}_m)^2
\]
гэж тооцно.

Мөн $i$-р хөрөнгийн регрессийн үлдэгдлийн
вариацыг
\[
\widehat{\sigma}_{\varepsilon,i}^{2}
=
\frac{1}{n}
\sum_{t=1}^{n}
\widehat{\varepsilon}_{ti}^{\,2}
\]
гэж тооцно.

Иймээс нэг хүчин зүйлийн зорилтот матриц
$T_{\mathrm{SF}}=(t_{ij})$-ын элементүүдийг
\[
t_{ij}
=
\begin{cases}
	\widehat{\beta}_i^2\widehat{\sigma}_m^2
	+
	\widehat{\sigma}_{\varepsilon,i}^{2},
	& i=j,\\[0.4em]
	\widehat{\beta}_i\widehat{\beta}_j
	\widehat{\sigma}_m^2,
	& i\neq j
\end{cases}
\]
гэж тодорхойлно \citep{ledoitwolf2003}.

Нэг хүчин зүйлийн зорилтот матрицыг матриц
хэлбэрээр
\[
T_{\mathrm{SF}}
=
\widehat{\boldsymbol{\beta}}
\widehat{\boldsymbol{\beta}}^{\top}
\widehat{\sigma}_m^2
+
\widehat D_{\varepsilon}
\]
гэж илэрхийлж болно.

Энд
\[
\widehat{\boldsymbol{\beta}}
=
(\widehat{\beta}_1,
\widehat{\beta}_2,
\ldots,
\widehat{\beta}_N)^{\top}
\]
нь регрессийн коэффициентүүдийн вектор, харин
\[
\widehat D_{\varepsilon}
=
\operatorname{diag}
\left(
\widehat{\sigma}_{\varepsilon,1}^{2},
\ldots,
\widehat{\sigma}_{\varepsilon,N}^{2}
\right)
\]
нь үлдэгдлийн вариацуудаас бүрдэх диагональ
матриц юм.

Нэг хүчин зүйлийн зорилтот матрицад суурилсан
шугаман агшилтын үнэлэгчийг
\[
\widehat{\Sigma}_{\mathrm{SF}}
=
(1-\widehat{\alpha}_{\mathrm{SF}})S
+
\widehat{\alpha}_{\mathrm{SF}}T_{\mathrm{SF}},
\qquad
0\leq\widehat{\alpha}_{\mathrm{SF}}\leq1,
\]
гэж тодорхойлно.

Энд $\widehat{\alpha}_{\mathrm{SF}}$ нь нэг хүчин
зүйлийн зорилтот матрицад харгалзах агшилтын
эрчим юм.

\paragraph{Агшилтын эрчмийн үнэлгээ.}

\citet{ledoitwolf2003} нь нэг хүчин зүйлийн
зорилтот матрицад суурилсан агшилтын эрчмийг
үнэлэхдээ түүврийн ковариацын матрицын
элементүүдийн түүвэрлэлтийн хэлбэлзэл,
түүврийн ковариацын матриц болон зорилтот
матрицын үнэлгээний алдаануудын хам хэлбэлзэл,
мөн эдгээр матрицуудын хоорондын зайг
харгалзан үзсэн.

Эхлээд түүврийн ковариацын матриц $S$ болон
нэг хүчин зүйлийн зорилтот матриц
$T_{\mathrm{SF}}$ хоорондын нормчилсон
квадрат Фробениусын зайг
\[
\widehat d_{\mathrm{SF}}
=
\left\|S-T_{\mathrm{SF}}\right\|^2=
\frac{1}{N}
\sum_{i=1}^{N}
\sum_{j=1}^{N}
(s_{ij}-t_{ij})^2
\]
гэж тодорхойлно.

Дараа нь түүврийн ковариацын матрицын
элементүүдийн түүвэрлэлтийн хэлбэлзлийг
нэгтгэн хэмжих хэмжигдэхүүнийг
\[
\widehat b_{\mathrm{SF}}
=
\frac{1}{nN}
\sum_{i=1}^{N}
\sum_{j=1}^{N}
\left[
\frac{1}{n}
\sum_{t=1}^{n}
\left\{
(y_{ti}-\overline{Y}_i)
(y_{tj}-\overline{Y}_j)
-s_{ij}
\right\}^{2}
\right]
\]
гэж тооцно.

Нэг хүчин зүйлийн зорилтот матриц нь
түүврийн өгөгдлөөс үнэлэгддэг учраас
$S$ болон $T_{\mathrm{SF}}$ матрицуудын
үнэлгээний алдаанууд хоорондоо
хамааралтай байж болно.

Иймээс агшилтын эрчмийг тооцоход
эдгээр үнэлгээний алдаануудын
хам хэлбэлзлийг харгалзах шаардлагатай.

Энэ хам хэлбэлзлийн нөлөөг
$\widehat c_{\mathrm{SF}}$ хэмжигдэхүүнээр
илэрхийлнэ.

Нэг хүчин зүйлийн зорилтот матрицын диагональ
элементүүд нь түүврийн ковариацын матрицын
диагональ элементүүдтэй тэнцүү буюу
$t_{ii}=s_{ii}$ байна.

\citet{ledoitwolf2003}-ийн Lemma 2-т
үндэслэн үнэлгээний алдаануудын
хам хэлбэлзлийг дараах байдлаар тооцно.

$i=j$ үед
\[
\widehat q_{ii}
=
\frac{1}{n}
\sum_{t=1}^{n}
\left\{
(y_{ti}-\overline{Y}_i)^2-s_{ii}
\right\}^{2}.
\]

Харин $i\neq j$ үед
\[
\widehat q_{ij}
=
\frac{1}{n}
\sum_{t=1}^{n}
\Big[
\widehat{\beta}_j
(y_{ti}-\overline{Y}_i)
(R_{m,t}-\overline{R}_m)
+
\widehat{\beta}_i
(y_{tj}-\overline{Y}_j)
(R_{m,t}-\overline{R}_m)
-
\widehat{\beta}_i\widehat{\beta}_j
(R_{m,t}-\overline{R}_m)^2
-t_{ij}
\Big]\times
\]
\[\times
\Big[
(y_{ti}-\overline{Y}_i)
(y_{tj}-\overline{Y}_j)-s_{ij}
\Big].
\]

Иймээс хам хэлбэлзлийн нийлбэрийн
нормчилсон үнэлгээг
\[
\widehat c_{\mathrm{SF}}
=
\frac{1}{nN}
\sum_{i=1}^{N}
\sum_{j=1}^{N}
\widehat q_{ij}
\]
гэж тодорхойлно.
Ингэснээр нэг хүчин зүйлийн зорилтот
матрицад харгалзах агшилтын эрчмийг
\[
\widehat{\alpha}_{\mathrm{SF}}
=
\max
\left\{
0,
\min
\left\{
1,
\frac{
	\widehat b_{\mathrm{SF}}
	-
	\widehat c_{\mathrm{SF}}
}{
	\widehat d_{\mathrm{SF}}
}
\right\}
\right\}
\]
гэж үнэлнэ.

Энэ нь \citet{ledoitwolf2003}-ийн
асимптотик оновчтой агшилтын эрчмийн
үнэлгээг нормчилсон квадрат
Фробениусын зайд нийцүүлэн
илэрхийлсэн хэлбэр юм.

Энд $\widehat d_{\mathrm{SF}}>0$
байхыг шаардана.

Эцэст нь нэг хүчин зүйлийн зорилтот
матрицад суурилсан агшилтын ковариацын
матрицыг
\[
\widehat{\Sigma}_{\mathrm{SF}}
=
(1-\widehat{\alpha}_{\mathrm{SF}})S
+
\widehat{\alpha}_{\mathrm{SF}}T_{\mathrm{SF}}
\]
гэж олно.

\paragraph{Тооцоолох алгоритм.}

Нэг хүчин зүйлийн зорилтот матрицад
суурилсан агшилтын ковариацын
матрицыг дараах үе шаттайгаар тооцно.

\begin{enumerate}
	
	\item Зах зээлийн дундаж өгөөжийг тооцно.
	\[
	\overline{R}_m
	=
	\frac{1}{n}
	\sum_{t=1}^{n}R_{m,t}.
	\]
	
	
	
	\item Зах зээлийн өгөөжийн түүврийн
	вариацыг тооцно.
	\[
	\widehat{\sigma}_m^2
	=
	\frac{1}{n}
	\sum_{t=1}^{n}
	(R_{m,t}-\overline{R}_m)^2.
	\]
	
	\item Хөрөнгө тус бүрийн регрессийн
	коэффициентийг тооцно.
	\[
	\widehat{\beta}_i
	=
	\frac{
		\sum_{t=1}^{n}
		(y_{ti}-\overline{Y}_i)
		(R_{m,t}-\overline{R}_m)
	}{
		\sum_{t=1}^{n}
		(R_{m,t}-\overline{R}_m)^2
	}.
	\]
	
	\item Регрессийн тогтмол коэффициент
	болон үлдэгдлийг тооцно.
	\[
	\widehat a_i
	=
	\overline{Y}_i
	-
	\widehat{\beta}_i\overline{R}_m,
	\]
	\[
	\widehat{\varepsilon}_{ti}
	=
	y_{ti}
	-
	\widehat a_i
	-
	\widehat{\beta}_iR_{m,t}.
	\]
	
	\item Хөрөнгө тус бүрийн үлдэгдлийн
	вариацыг тооцно.
	\[
	\widehat{\sigma}_{\varepsilon,i}^{2}
	=
	\frac{1}{n}
	\sum_{t=1}^{n}
	\widehat{\varepsilon}_{ti}^{\,2}.
	\]
	
	\item Нэг хүчин зүйлийн зорилтот
	матрицыг байгуулна.
	\[
	t_{ij}
	=
	\begin{cases}
		\widehat{\beta}_i^2
		\widehat{\sigma}_m^2
		+
		\widehat{\sigma}_{\varepsilon,i}^{2},
		& i=j,\\[0.4em]
		\widehat{\beta}_i\widehat{\beta}_j
		\widehat{\sigma}_m^2,
		& i\neq j.
	\end{cases}
	\]
	
	Эдгээр элементүүдээр
	$T_{\mathrm{SF}}=(t_{ij})$
	матрицыг байгуулна.
	
	\item Түүврийн ковариацын матриц болон
	зорилтот матрицын хоорондын нормчилсон
	квадрат Фробениусын зайг тооцно.
	\[
	\widehat d_{\mathrm{SF}}
	=
	\frac{1}{N}
	\sum_{i=1}^{N}
	\sum_{j=1}^{N}
	(s_{ij}-t_{ij})^2.
	\]
	
	\item Түүврийн ковариацын матрицын
	элементүүдийн түүвэрлэлтийн
	хэлбэлзлийг тооцно.
	\[
	\widehat b_{\mathrm{SF}}
	=
	\frac{1}{nN}
	\sum_{i=1}^{N}
	\sum_{j=1}^{N}
	\left[
	\frac{1}{n}
	\sum_{t=1}^{n}
	\left\{
	(y_{ti}-\overline{Y}_i)
	(y_{tj}-\overline{Y}_j)
	-s_{ij}
	\right\}^{2}
	\right].
	\]
	
	\item Түүврийн ковариацын матриц
	болон нэг хүчин зүйлийн зорилтот
	матрицын үнэлгээний алдаануудын
	хам хэлбэлзлийг тооцно.
	
	$i=j$ үед
	\[
	\widehat q_{ii}
	=
	\frac{1}{n}
	\sum_{t=1}^{n}
	\left\{
	(y_{ti}-\overline{Y}_i)^2-s_{ii}
	\right\}^{2}.
	\]
	
	$i\neq j$ үед
	\[
	\widehat q_{ij}
	=
	\frac{1}{n}
	\sum_{t=1}^{n}
	\Big[
	\widehat{\beta}_j
	(y_{ti}-\overline{Y}_i)
	(R_{m,t}-\overline{R}_m)
	+
	\widehat{\beta}_i
	(y_{tj}-\overline{Y}_j)
	(R_{m,t}-\overline{R}_m)
	-
	\widehat{\beta}_i\widehat{\beta}_j
	(R_{m,t}-\overline{R}_m)^2
	-t_{ij}
	\Big]\times
	\]
	\[\times
	\Big[
	(y_{ti}-\overline{Y}_i)
	(y_{tj}-\overline{Y}_j)-s_{ij}
	\Big].
	\]
	
	Дараа нь
	\[
	\widehat c_{\mathrm{SF}}
	=
	\frac{1}{nN}
	\sum_{i=1}^{N}
	\sum_{j=1}^{N}
	\widehat q_{ij}
	\]
	гэж тооцно.
	
	\item Агшилтын эрчмийг тооцно.
	
	$\widehat d_{\mathrm{SF}}>0$ үед
	\[
	\widehat{\alpha}_{\mathrm{SF}}
	=
	\max
	\left\{
	0,
	\min
	\left\{
	1,
	\frac{
		\widehat b_{\mathrm{SF}}
		-
		\widehat c_{\mathrm{SF}}
	}{
		\widehat d_{\mathrm{SF}}
	}
	\right\}
	\right\}.
	\]
	
	\item Агшилтын ковариацын
	матрицыг байгуулна.
	\[
	\widehat{\Sigma}_{\mathrm{SF}}
	=
	(1-\widehat{\alpha}_{\mathrm{SF}})S
	+
	\widehat{\alpha}_{\mathrm{SF}}T_{\mathrm{SF}}.
	\]
	
\end{enumerate}

% --------------------------------------------------
\subsection{Шугаман бус агшилтын үнэлгээ}
% --------------------------------------------------

Шугаман агшилтын аргын хязгаарлалтыг даван туулах зорилгоор шугаман бус агшилтын (Nonlinear Shrinkage) аргуудыг боловсруулсан. Эдгээр аргууд нь түүврийн ковариацын матрицын спектрийн задлалд суурилдаг бөгөөд урьдчилан тодорхойлсон нэг зорилтот матриц руу бүхэлд нь агшаахын оронд түүврийн ковариацын матрицын хувийн утгуудыг ялгаатай хэмжээгээр засварладаг.

Түүврийн ковариацын матрицын спектрийн задлалыг
\[
S
=
U\Lambda U^{\top}
\]
гэж тэмдэглэе. Энд $U=(u_1,u_2,\ldots,u_N)$ нь $S$ матрицын хувийн векторуудаас бүрдэх ортогональ матриц, харин
\[
\Lambda
=
\operatorname{diag}
(\lambda_1,\lambda_2,\ldots,\lambda_N)
\]
нь түүврийн ковариацын матрицын хувийн утгуудыг агуулсан диагональ матриц юм.

Шугаман бус агшилтын үнэлэгчийг
\[
\widehat{\Sigma}_{\mathrm{NLS}}
=
U\widehat{D}U^{\top}
\]
гэж тодорхойлно. Энд
\[
\widehat{D}
=
\operatorname{diag}
(\widehat{d}_1,\widehat{d}_2,\ldots,\widehat{d}_N)
\]
нь агшааж засварласан хувийн утгуудаас бүрдэх диагональ матриц бөгөөд $\widehat{d}_i$ нь $\lambda_i$ хувийн утгад харгалзах засварласан хувийн утга юм.

Шугаман бус агшилтын арга нь түүврийн ковариацын матрицын хувийн векторуудыг хадгалж, харин хувийн утга бүрийг өгөгдөлд тулгуурлан ялгаатай хэмжээгээр засварладаг. Иймээс хувийн утгуудын хэмжээ болон тэдгээрийн спектрийн тархалт дахь байрлалаас хамааран агшилтын хэмжээ өөр өөр байж болно.

Шугаман агшилтын аргын хувьд түүврийн ковариацын матриц болон зорилтот матрицыг нэг агшилтын эрчмээр хослуулдаг бол шугаман бус агшилтын арга нь хувийн утга бүрд тусгайлан тохируулсан хувиргалт хийдгээрээ ялгаатай. Ингэснээр түүврийн ковариацын матрицын хувийн утгуудын үнэлгээний алдааг бууруулах боломжийг олгодог.

Ledoit болон Wolf (2012) нар шугаман бус агшилтын
аргын онолын үндэслэлийг боловсруулж, улмаар
Ledoit болон Wolf (2015) нар QuEST
(Quantized Eigenvalues Sampling Transform)
алгоритмд суурилсан тоон хэрэгжилтийг хөгжүүлсэн.
Энэхүү арга нь тоон оновчлол шаарддаг тул
тооцооллын хувьд харьцангуй төвөгтэй байдаг.

Үүний дараа \citet{ledoitwolf2020} санамсаргүй
матрицын онол болон спектрийн тархалтын
шинж чанарт тулгуурласан аналитик шугаман бус
агшилтын аргыг боловсруулсан. Уг арга нь
түүврийн ковариацын матрицын хувийн утгуудыг
аналитик илэрхийлэлд суурилан засварладаг.
Энэхүү судалгаанд аналитик шугаман бус
агшилтын аргыг ашиглах бөгөөд түүний
математик үндэслэл болон тооцоолох
алгоритмыг дараагийн хэсэгт авч үзнэ.




% --------------------------------------------------
\subsubsection{Аналитик шугаман бус агшилтын үнэлгээ}
% --------------------------------------------------

Аналитик шугаман бус агшилтын үнэлгээ
(Analytical Nonlinear Shrinkage Estimation) нь
санамсаргүй матрицын онол (Random Matrix Theory)-д
тулгуурлан түүврийн ковариацын матрицын хувийн
утгуудыг засварлах арга юм. Шугаман агшилтын
аргуудаас ялгаатай нь энэхүү арга нь хувийн
утга бүрийн хэмжээ болон спектрийн тархалт дахь
байрлалыг харгалзан ялгаатай агшилт хийдэг.

\citet{ledoitwolf2020} нь спектрийн тархалтын
нягт болон түүний Hilbert хувиргалтыг kernel
аргаар үнэлэх замаар аналитик шугаман бус
агшилтын аргыг боловсруулсан. Уг арга нь
түүврийн ковариацын матрицын хувийн
векторуудыг хадгалж, хувийн утгуудыг
засварлах замаар ковариацын матрицын
үнэлгээний алдааг бууруулахад чиглэдэг.

Судалгаанд хамрагдсан хөрөнгийн тоог $N$,
ажиглалтын тоог $n$ гэж тэмдэглэж,
тэдгээрийн харьцааг
\[
c=\frac{N}{n}
\]
гэж тодорхойлно. Аргын үндсэн онолын хувьд
$N,n\to\infty$ үед
\[
\frac{N}{n}\longrightarrow c\in(0,1)
\]
байх асимптотик нөхцөлийг авч үздэг.

Түүврийн ковариацын матрицын спектрийн
задлалыг
\[
S=U\Lambda U^{\top}
\]
гэж тэмдэглэе. Энд
\[
U=(u_1,u_2,\ldots,u_N),
\qquad
U^{\top}U=I_N,
\]
нь хувийн векторуудаас бүрдэх ортогональ
матриц, харин
\[
\Lambda
=
\operatorname{diag}
(\lambda_1,\lambda_2,\ldots,\lambda_N)
\]
нь хувийн утгуудыг агуулсан диагональ
матриц юм. Хувийн утгуудыг
\[
\lambda_1\geq\lambda_2\geq\cdots
\geq\lambda_N>0
\]
гэж эрэмбэлнэ.

Аналитик шугаман бус агшилтын үнэлэгчийг
\[
\widehat{\Sigma}_{\mathrm{NLS}}
=
U\widehat{D}U^{\top}
=
\sum_{i=1}^{N}
\widehat{d}_i u_i u_i^{\top}
\]
гэж тодорхойлно. Энд
\[
\widehat{D}
=
\operatorname{diag}
(\widehat{d}_1,\ldots,\widehat{d}_N)
\]
нь агшсан хувийн утгуудаас бүрдэх
диагональ матриц юм.

\paragraph{Агшсан хувийн утгуудын үнэлгээ.}

\citet{ledoitwolf2020}-ийн (4.3) томьёоны
дагуу агшсан хувийн утгуудыг
\[
\widehat{d}_i
=
\frac{\lambda_i}
{
	\left(
	1-c-c\pi\lambda_i
	\widehat{Hf}(\lambda_i)
	\right)^2
	+
	\left(
	c\pi\lambda_i
	\widehat{f}(\lambda_i)
	\right)^2
}
\]
гэж тооцно.

Энд $\widehat{f}(\lambda_i)$ нь
спектрийн тархалтын нягтын үнэлгээ,
$\widehat{Hf}(\lambda_i)$ нь уг нягтын
Hilbert хувиргалтын үнэлгээ юм.
Эдгээрийг түүврийн ковариацын матрицын
хувийн утгуудад тулгуурлан kernel
аргаар тооцдог.

Иймээс агшсан хувийн утга нь зөвхөн
тухайн $\lambda_i$-ийн хэмжээнээс бус,
спектрийн тархалтын нягт болон түүний
Hilbert хувиргалтын утгаас хамаарна.

\paragraph{Тооцоолох алгоритм.}

\citet{ledoitwolf2020}-ийн Section 4.7-д
үндэслэн аналитик шугаман бус агшилтын
ковариацын матрицыг дараах үе шаттайгаар
тооцно.

\begin{enumerate}
	
	
	\item Түүврийн ковариацын матрицын
	спектрийн задлалыг байгуулна.
	\[
	S=U\Lambda U^{\top}.
	\]
	
	\item Хөрөнгийн тоо болон ажиглалтын
	тооны харьцааг тооцно.
	\[
	c=\frac{N}{n}.
	\]
	
	\item Спектрийн нягтыг үнэлэхэд
	ашиглах bandwidth параметрийг
	\[
	h=n^{-1/3}
	\]
	гэж сонгоно. Хувийн утга бүрд
	харгалзах локал bandwidth нь
	\[
	h_j=\lambda_jh,
	\qquad j=1,\ldots,N
	\]
	байна.
	
	\item Спектрийн нягтыг Epanechnikov
	kernel ашиглан үнэлнэ.
	\[
	\widehat{f}(\lambda_i)
	=
	\frac{1}{N}
	\sum_{j=1}^{N}
	\frac{1}{h_j}
	k\left(
	\frac{\lambda_i-\lambda_j}{h_j}
	\right),
	\]
	энд
	\[
	k(x)
	=
	\begin{cases}
		\dfrac{3}{4\sqrt{5}}
		\left(1-\dfrac{x^2}{5}\right),
		& |x|\leq\sqrt{5},\\[0.5em]
		0,
		& |x|>\sqrt{5}.
	\end{cases}
	\]
	
	\item Спектрийн нягтын Hilbert
	хувиргалтыг үнэлнэ.
	\[
	\widehat{Hf}(\lambda_i)
	=
	\frac{1}{N}
	\sum_{j=1}^{N}
	\frac{1}{h_j}
	\ell\left(
	\frac{\lambda_i-\lambda_j}{h_j}
	\right),
	\]
	энд $\ell(x)=Hk(x)$ бөгөөд
	\[
	\ell(x)
	=
	\begin{cases}
		-\dfrac{3x}{10\pi}
		+
		\dfrac{3}{4\sqrt{5}\pi}
		\left(1-\dfrac{x^2}{5}\right)
		\ln\left|
		\dfrac{\sqrt{5}-x}{\sqrt{5}+x}
		\right|,
		& |x|\neq\sqrt{5},\\[1em]
		-\dfrac{3x}{10\pi},
		& |x|=\sqrt{5}.
	\end{cases}
	\]
	
	\item Агшсан хувийн утгуудыг
	тооцно.
	\[
	\widehat{d}_i
	=
	\frac{\lambda_i}
	{
		\left(
		1-c-c\pi\lambda_i
		\widehat{Hf}(\lambda_i)
		\right)^2
		+
		\left(
		c\pi\lambda_i
		\widehat{f}(\lambda_i)
		\right)^2
	}.
	\]
	
	\item Агшилтын ковариацын матрицыг
	байгуулна.
	\[
	\widehat{\Sigma}_{\mathrm{NLS}}
	=
	U
	\operatorname{diag}
	(\widehat{d}_1,\ldots,\widehat{d}_N)
	U^{\top}.
	\]
	
\end{enumerate}



% ----------------------------------------------------------
\subsection{Эмпирик үнэлгээний арга зүй}
\label{subsec:empirical_framework}
% ----------------------------------------------------------

Энэхүү судалгаанд түүврийн ковариацын матриц болон
агшилтын таван үнэлэгчийг ижил long-only GMVP
загварын хүрээнд түүврээс гадуурх
(out-of-sample, OOS) гүйцэтгэлээр харьцуулна.
Ковариацын матрицын үнэлгээний алдаа бага байх
нь багцын бодит эрсдэл заавал бага байна гэсэн
үг биш тул үнэлэгчдийн гүйцэтгэлийг багцын
түвшинд үнэлэх шаардлагатай.

Энэ зорилгоор
\citep{ledoitwolf2003,demiguel2009}
судалгаанд хэрэглэсэн эмпирик үнэлгээний
зарчимд тулгуурлан шилжих цонхны аргыг
ашиглана. Бүх үнэлэгчийн хувьд хөрөнгө
оруулалтын хязгаарлалт, түүврээс гадуурх
үнэлгээний хугацаа болон гүйцэтгэлийн
шалгуурыг ижил байлгана.

% ----------------------------------------------------------
% 2.6.1. Rolling-Window Out-of-Sample Evaluation
% ----------------------------------------------------------

\subsubsection{Шилжих цонхны аргаар түүврээс гадуурх
	үнэлгээ хийх}
\label{subsubsec:rolling_window}

Шилжих цонхны (rolling-window) арга нь тогтмол
урттай түүхэн өгөгдлөөр ковариацын матрицыг
үнэлж, багцын жинг тодорхойлсны дараа уг
жинг дараагийн сарын өгөөжид хэрэглэнэ.

Нийт сарын өгөөжийн ажиглалтын тоог $n$,
хөрөнгийн тоог $N$, сургалтын цонхны
уртыг $T$ гэж тэмдэглэе.
$t$ сарын өгөөжийн вектор нь
\[
\boldsymbol{R}_t
=
(R_{1,t},R_{2,t},\ldots,R_{N,t})^\top.
\]

$t$ сарын эцэст ашиглах сургалтын өгөгдлийг
\[
\mathcal{D}_t
=
\{\boldsymbol{R}_{t-T+1},\ldots,
\boldsymbol{R}_t\},
\qquad t=T,\ldots,n-1
\]
гэж тодорхойлно.

Ковариацын матрицын үнэлэгчийг
$k=0,1,\ldots,5$ гэж тэмдэглэе.
Энд $k=0$ нь түүврийн ковариацын матриц,
$k=1,\ldots,5$ нь агшилтын үнэлэгчид юм.
Сургалтын цонх бүрд ковариацын матрицыг
\[
\widehat{\boldsymbol{\Sigma}}_t^{(k)}
=
\widehat{\boldsymbol{\Sigma}}^{(k)}
(\mathcal{D}_t)
\]
гэж үнэлнэ. Нэг хүчин зүйлийн үнэлэгчийн
хувьд сургалтын цонхтой ижил хугацааны
зах зээлийн индексийн өгөөжийг нэмж ашиглана.

Үнэлсэн ковариацын матрицад үндэслэн
багцын оновчтой жинг
\[
\begin{aligned}
	\widehat{\boldsymbol{w}}_t^{(k)}
	&\in
	\underset{\boldsymbol{w}}{\arg\min}
	\quad
	\boldsymbol{w}^{\top}
	\widehat{\boldsymbol{\Sigma}}_t^{(k)}
	\boldsymbol{w},\\
	\text{subject to}\quad
	&\boldsymbol{1}^{\top}\boldsymbol{w}=1,\\
	&w_i\geq 0,\qquad i=1,\ldots,N
\end{aligned}
\]
гэсэн long-only GMVP бодлогоор тодорхойлно.

Тэгвэл дараагийн сарын багцын өгөөж нь
\[
r_{p,t+1}^{(k)}
=
\left(\widehat{\boldsymbol{w}}_t^{(k)}\right)^\top
\boldsymbol{R}_{t+1}
\]
байна.

Сургалтын цонхыг сар бүр нэг алхмаар
урагшлуулж, ковариацын матриц, агшилтын
эрчим болон багцын жинг дахин тооцно.
Ирээдүйн өгөөжийг параметрийн үнэлгээнд
ашиглахгүй.

Үндсэн шинжилгээнд $N=10$, $n=172$,
$T=120$ тул түүврээс гадуурх ажиглалтын тоо
\[
M=n-T=172-120=52
\]
байна. Ийнхүү үнэлэгч бүрийн хувьд
ижил 52 сарын OOS өгөөжийн цуваа үүсгэнэ.

% ----------------------------------------------------------
% 2.6.2. Portfolio Performance Measures
% ----------------------------------------------------------

\subsubsection{Багцын гүйцэтгэлийн үзүүлэлтүүд}
\label{subsubsec:performance_measures}

Ковариацын үнэлэгчдийн гүйцэтгэлийг
OOS хэлбэлзэл, харьцангуй дисперсийн
бууралт (RVR) болон Sharpe-ийн
харьцаагаар үнэлнэ.

\paragraph{Түүврээс гадуурх хэлбэлзэл
	(Out-of-Sample Volatility).}

$k$-р үнэлэгчээр байгуулсан багцын
дундаж сарын өгөөж болон түүврийн
вариацыг
\[
\overline{r}_p^{(k)}
=
\frac{1}{M}
\sum_{t=T+1}^{n}r_{p,t}^{(k)},
\]
\[
\widehat{\sigma}_{p,k}^{\,2}
=
\frac{1}{M-1}
\sum_{t=T+1}^{n}
\left(
r_{p,t}^{(k)}
-
\overline{r}_p^{(k)}
\right)^2
\]
гэж тодорхойлно.

Тэгвэл сарын болон жилийн хэлбэлзэл нь
\[
\widehat{\sigma}_{p,k}
=
\sqrt{\widehat{\sigma}_{p,k}^{\,2}},
\qquad
\widehat{\sigma}_{p,k}^{\,ann}
=
\sqrt{12}\,\widehat{\sigma}_{p,k}
\]
байна. Жилийн хэлбэлзэлд шилжүүлэхдээ
сарын өгөөжүүдийн автоковариацыг тэг
гэж үзнэ.

\paragraph{Харьцангуй дисперсийн бууралт
	(Relative Variance Reduction).}

Агшилтын үнэлэгчийн OOS вариацыг
түүврийн ковариацын матрицад суурилсан
GMVP-тэй харьцуулж,
\[
RVR_k^{(\%)}
=
100\left(
1-
\frac{\widehat{\sigma}_{p,k}^{\,2}}
{\widehat{\sigma}_{p,0}^{\,2}}
\right),
\qquad k=1,\ldots,5
\]
гэж тооцно.

$RVR_k^{(\%)}>0$ бол агшилтын үнэлэгчийн
багц суурь загвараас бага вариацтай,
харин $RVR_k^{(\%)}<0$ бол их вариацтай
байна.

\paragraph{Sharpe-ийн харьцаа
	(Sharpe Ratio).}

Сарын эрсдэлгүй өгөөжийг $r_{f,t}$
гэвэл илүүдэл өгөөж нь
\[
r_{e,t}^{(k)}
=
r_{p,t}^{(k)}-r_{f,t}
\]
байна. Сарын Sharpe-ийн харьцааг
\[
\widehat{SR}_k
=
\frac{\overline{r}_e^{(k)}}
{\widehat{\sigma}_{e,k}}
\]
гэж тодорхойлно.

Энэхүү судалгаанд $r_{f,t}=0$
гэж авч үзсэн тул
\[
\widehat{SR}_k
=
\frac{\overline{r}_p^{(k)}}
{\widehat{\sigma}_{p,k}},
\qquad
\widehat{SR}_k^{\,ann}
=
\sqrt{12}\,\widehat{SR}_k.
\]

Жилийн Sharpe-ийн харьцаанд шилжүүлэхдээ
мөн сарын өгөөжийн автоковариацыг тэг
гэж үзнэ. GMVP-ийн зорилго нь вариацыг
хамгийн бага байлгах тул Sharpe-ийн
харьцааг нэмэлт гүйцэтгэлийн үзүүлэлт
болгон ашиглана.

% ----------------------------------------------------------
% 2.6.3. Portfolio Weight Stability and Turnover
% ----------------------------------------------------------

\subsubsection{Багцын жингийн тогтвортой
	байдал ба эргэлт}
\label{subsubsec:weight_stability}

Ковариацын үнэлэгчийн сонголт нь
багцын жингийн бүтэц болон дахин
тэнцвэржүүлэлтийн хэмжээнд нөлөөлж
болох тул багцын эргэлт, жингийн
төвлөрлийг үнэлнэ.


\paragraph{Багцын эргэлт (Portfolio Turnover).}

$t$ сард ашиглах багцын жинг
$\widehat{\boldsymbol{w}}_{t-1}^{(k)}$
гэж тэмдэглэе. Тухайн сарын өгөөж
хэрэгжсэний дараах, дахин
тэнцвэржүүлэхээс өмнөх жин нь
\[
\widetilde{w}_{i,t}^{(k)}
=
\frac{
	\widehat{w}_{i,t-1}^{(k)}(1+R_{i,t})
}{
	1+r_{p,t}^{(k)}
}
\]
байна.

Дараагийн сард ашиглах шинэ
оновчтой жинд шилжихэд шаардагдах
багцын эргэлтийг
\[
TO_t^{(k)}
=
\frac{1}{2}
\sum_{i=1}^{N}
\left|
\widehat{w}_{i,t}^{(k)}
-
\widetilde{w}_{i,t}^{(k)}
\right|
\]
гэж тодорхойлно.

Түүврээс гадуурх дундаж эргэлт нь
\[
\overline{TO}^{(k)}
=
\frac{1}{M-1}
\sum_{t=T+1}^{n-1}TO_t^{(k)}
\]
байна.

Ийнхүү 52 сарын OOS хугацаанд
51 удаагийн дахин тэнцвэржүүлэлтийг
тооцно. Эхний багц бүрдүүлэлтийг
эргэлтийн тооцоонд оруулахгүй.
Энд тооцсон үндсэн OOS өгөөж нь арилжааны
зардлын өмнөх өгөөж юм. Арилжааны зардлын
нөлөөг нэмэлт шинжилгээнд тусад нь үнэлнэ.


\paragraph{Багцын жингийн төвлөрөл
	(Weight Concentration).}

Жингийн төвлөрлийг
Herfindahl--Hirschman Index (HHI)-ээр
\[
HHI_t^{(k)}
=
\sum_{i=1}^{N}
\left(\widehat{w}_{i,t}^{(k)}\right)^2
\]
гэж хэмжинэ.

Long-only GMVP-ийн хувьд
\[
\frac{1}{N}
\leq
HHI_t^{(k)}
\leq 1.
\]

Түүврээс гадуурх дундаж төвлөрөл нь
\[
\overline{HHI}^{(k)}
=
\frac{1}{M}
\sum_{t=T}^{n-1}HHI_t^{(k)}
\]
байна. $HHI$ өндөр байх нь багцын
жин цөөн хөрөнгөд төвлөрснийг
илэрхийлнэ.

% ----------------------------------------------------------
% 2.6.4. Statistical Comparison
% ----------------------------------------------------------

\subsubsection{Ковариацын үнэлэгчдийн
	статистик харьцуулалт}
\label{subsubsec:statistical_comparison}

Агшилтын үнэлэгчдийн OOS гүйцэтгэлийн
ялгааг түүврийн ковариацын матрицад
суурилсан GMVP-тэй харьцуулан
статистикаар үнэлнэ.

$k$-р агшилтын үнэлэгч болон суурь
загварын сарын квадрат өгөөжийн
ялгааг
\[
d_t^{(k)}
=
\left(r_{p,t}^{(k)}\right)^2
-
\left(r_{p,t}^{(0)}\right)^2,
\qquad t=T+1,\ldots,n
\]
гэж тодорхойлно.

Дараах хоёр талт таамаглалыг шалгана:
\[
H_0:E[d_t^{(k)}]=0,
\qquad
H_1:E[d_t^{(k)}]\neq 0.
\]

Дундаж ялгааны стандарт алдааг
Newey--West төрлийн HAC аргаар,
гурван лаг ашиглан үнэлнэ.
Шалгах статистик нь
\[
t_{HAC}^{(k)}
=
\frac{\overline{d}^{(k)}}
{\widehat{SE}_{HAC}
	(\overline{d}^{(k)})},
\qquad
\overline{d}^{(k)}
=
\frac{1}{M}
\sum_{t=T+1}^{n}d_t^{(k)}
\]
байна.

Хоёр талт $p$-утгыг асимптотик
нормаль ойролцооллоор тооцож,
таван агшилтын үнэлэгчийг суурь
загвартай харьцуулсан олон шалгалтад
Holm-ийн засвар хэрэглэнэ.

Квадрат өгөөжийн дундаж ялгаа нь
дисперсийн ялгаатай заавал тэнцүү биш.
Учир нь
\[
E[r_p^2]
=
\operatorname{Var}(r_p)
+
\left(E[r_p]\right)^2.
\]

Иймээс HAC шалгалтыг RVR болон
OOS хэлбэлзлийн үр дүнтэй хамтатган
тайлбарлана. Мөн OOS хугацаа
52 сар тул статистик шалгалтын
чадал хязгаарлагдмал байж болно.

% ----------------------------------------------------------
% 2.6.5. Sensitivity to N/T
% ----------------------------------------------------------

\subsubsection{Хөрөнгийн тоо ба сургалтын
	ажиглалтын тооны харьцааны мэдрэмжийн шинжилгээ}
\label{subsubsec:nt_sensitivity}

Ковариацын матрицын үнэлгээний алдаа
нь хөрөнгийн тоо $N$ болон сургалтын
ажиглалтын тоо $T$-ийн харьцаанаас
хамаарч болно. Уг харьцааг
\[
c=\frac{N}{T}
\]
гэж тодорхойлно.

Энэхүү судалгаанд $N=10$ гэж
тогтмол байлгаж,
\[
T\in\{60,90,120\}
\]
гэсэн сургалтын цонхуудыг авч үзнэ.
Харгалзах $N/T$ харьцаа нь
\[
\frac{10}{60}\approx0.1667,\qquad
\frac{10}{90}\approx0.1111,\qquad
\frac{10}{120}\approx0.0833
\]
байна.

Сургалтын цонхны уртаас үүдэх
OOS хугацааны ялгааг арилгахын
тулд бүх хувилбарт хамгийн урт
сургалтын цонхонд харгалзах
ижил 52 сарын OOS хугацааг ашиглана.

Цонхны урт бүрд зургаан үнэлэгчийг
дахин тооцож, багцын хэлбэлзэл,
RVR, Sharpe-ийн харьцаа, эргэлт
болон төвлөрлийг харьцуулна.

Энэ шинжилгээнд $N$ тогтмол тул
$N/T$ харьцааны өөрчлөлт нь сургалтын
цонхны уртын өөрчлөлттэй давхцана.
Иймээс тэдгээрийн нөлөөг бүрэн
тусгаарлан тайлбарлахгүй.

% ----------------------------------------------------------
% 2.6.6. Portfolio Selection Robustness
% ----------------------------------------------------------

\subsubsection{Багцын сонголтын
	robustness шинжилгээ}
\label{subsubsec:portfolio_robustness}

Агшилтын үнэлэгчдийн гүйцэтгэл
сонгосон хөрөнгүүдийн бүрэлдэхүүнээс
хамаарах эсэхийг үнэлэхийн тулд
корреляцийн ялгаатай бүтэцтэй
150 багцад шинжилгээг давтан
гүйцэтгэнэ.

Сарын өгөөжийн бүрэн өгөгдөлтэй
337 хувьцаанаас тус бүр 10 хөрөнгөтэй
багцуудыг сонгоно. Энэхүү 337 хувьцааны олонлог нь
2012--2026 оны хугацаанд бүрэн өгөгдөлтэй байх
шалгуураар бүрдсэн тул үр дүнг уг сонгогдсон
олонлогийн хүрээнд тайлбарлана. Өгөгдлийн бүрэн
байдлын шалгуураас үүдэх сонголтын хазайлтыг
энэ шинжилгээ бүрэн арилгахгүй. Корреляцыг
эхний 120 сарын сургалтын өгөөжөөр
тооцож, багцын дундаж хос
корреляцыг
\[
\overline{\rho}
=
\frac{2}{N(N-1)}
\sum_{i<j}\rho_{ij}
\]
гэж тодорхойлно.

Багцуудыг өндөр (High),
дунд (Medium), бага (Low)
корреляцтай гурван ангиллаар
авч үзнэ. Үндсэн шинжилгээнд сонгосон High, Medium,
Low багцуудын эхний 120 сарын өгөөжөөр
тооцсон дундаж хос корреляц нь харгалзан
$0.557048$, $0.303609$, $0.101523$ байна.
Robustness шинжилгээнд эдгээр утгыг
зорилтот корреляц болгон ашиглана.

Багц сонгохдоо 337 хувьцаанаас
10 хувьцааг буцаалтгүйгээр
санамсаргүй сонгох туршилтыг
250000 удаа гүйцэтгэж,
давхардсан нэр дэвшигч
багцуудыг хасна.

Үлдсэн багцуудыг зорилтот
корреляцаас хазайх абсолют
зөрүүгээр эрэмбэлж, ангилал
бүрд 50 багц сонгоно.
Нэг ангиллын сонгогдсон
дурын хоёр багцын нийтлэг
хувьцааны тоо зургаагаас
хэтрэхгүй байх хязгаарлалт
тавина.

Ийнхүү сонгосон 150 багц бүрд
зургаан үнэлэгчээр long-only
GMVP байгуулж, 120 сарын
шилжих цонх болон ижил
52 сарын OOS хугацаагаар
үнэлнэ. Багц сонгох үе шатанд
OOS өгөөжийг ашиглахгүй.

Багц тус бүрийн RVR-ийг тооцож,
ангилал бүрийн дундаж, медиан
болон эерэг RVR-тэй багцын
эзлэх хувийг харьцуулна.

Эдгээр 150 багц нь бүх боломжит
10 хувьцаатай багцуудаас энгийн
санамсаргүй түүврийн аргаар
сонгогдоогүй, харин зорилтот
корреляцийн түвшинд ойр байх
нөхцөлөөр сонгогдсон болохыг
үр дүнг тайлбарлахдаа харгалзан үзнэ.

% ----------------------------------------------------------
% 2.6.7. Block Bootstrap Confidence Intervals
% ----------------------------------------------------------

\subsubsection{RVR-ийн bootstrap
	итгэх завсар}
\label{subsubsec:bootstrap_rvr}

150 багцын robustness шинжилгээнд
ажиглагдсан дундаж RVR-ийн
тодорхой бус байдлыг үнэлэхийн
тулд circular moving-block
bootstrap аргыг ашиглана.

Корреляцын ангиллыг
$g\in\{H,M,L\}$,
тухайн ангиллын багцын
тоог $P_g=50$ гэж тэмдэглэе.
$j$-р багцын $k$-р үнэлэгчийн
RVR-ийг $RVR_{g,j,k}^{(\%)}$
гэвэл ангиллын дундаж RVR нь
\[
\overline{RVR}_{g,k}^{(\%)}
=
\frac{1}{P_g}
\sum_{j=1}^{P_g}
RVR_{g,j,k}^{(\%)}
\]
байна.

Bootstrap давталт бүрд 52 сарын
OOS хугацаанаас зургаан сарын
урттай блокуудын эхлэх байрлалыг
санамсаргүй сонгоно. Блокийг
цувааны төгсгөлд хүрвэл эхлэлээс
үргэлжлүүлэх circular зарчмаар
байгуулж, сонгосон блокуудыг
нэгтгэн 52 ажиглалттай түүвэр
үүсгэнэ.

Тухайн давталтад сонгосон
сарын индексүүдийг бүх
150 багц болон зургаан
үнэлэгчид ижил хэрэглэнэ.
Ингэснээр багцуудын өгөөжийн
хоорондын хамаарлыг дахин
түүвэрлэх явцад хадгална.

Bootstrap түүвэр бүрд багц
тус бүрийн OOS вариац,
RVR болон ангиллын дундаж
RVR-ийг шинээр тооцно.

Нийт $B=3000$ давталт
гүйцэтгэж, давталтуудаас
гарсан дундаж RVR-ийн
тархалтын $2.5$ болон
$97.5$ хувийн квантилуудаар
95\%-ийн percentile bootstrap
итгэх завсрыг байгуулна.

Итгэх завсар тэгийг агуулаагүй
тохиолдолд тухайн ангиллын
дундаж RVR тэгээс ялгаатай
байж болохыг илтгэх боловч
энэ нь олон харьцуулалтын
засвартай шалгалттай адил биш. Эдгээр нь олон харьцуулалтын
засвар хийгдээгүй итгэх завсрууд
тул үр дүнг болгоомжтой тайлбарлана.

Харин итгэх завсар тэгийг
агуулсан тохиолдолд эерэг
эсвэл сөрөг нөлөөг статистикийн
хувьд баталгаатай гэж
дүгнэхгүй.

Энэхүү bootstrap шинжилгээнд
сонгосон 150 багцын
бүрэлдэхүүнийг тогтмол
байлгаж, зөвхөн OOS
хугацааны ажиглалтуудыг
дахин түүвэрлэнэ.

Иймээс итгэх завсар нь
тухайн сонгосон багцуудын
OOS гүйцэтгэлийн тодорхой
бус байдлыг илэрхийлэх
бөгөөд бүх боломжит
хувьцааны багцад шууд
ерөнхийлөх үндэслэл
болохгүй.

% ----------------------------------------------------------
% 2.6.8. Equal-Weight Benchmark and Transaction Costs
% ----------------------------------------------------------

\subsubsection{Тэнцүү жинтэй багц ба арилжааны зардлын шинжилгээ}
\label{subsubsec:equal_weight_costs}

GMVP загваруудын практик гүйцэтгэлийг
харьцуулах нэмэлт жишиг болгон хөрөнгө бүрд
ижил жин оноодог $1/N$ багцыг ашиглана.
Тухайн багцын зорилтот жин нь
\[
\boldsymbol{w}^{EW}=\frac{1}{N}\boldsymbol{1},
\qquad N=10,
\]
бөгөөд сарын эхэнд дахин тэнцвэржүүлж,
зорилтот жинг хадгална. Иймээс $t$ сарын
жишиг багцын зардлын өмнөх өгөөж нь
\[
r_{p,t}^{EW}
=\left(\boldsymbol{w}^{EW}\right)^\top
\boldsymbol{R}_t
=\frac{1}{N}\sum_{i=1}^{N}R_{i,t}.
\]

$1/N$ багц болон зургаан GMVP загварт
OOS хугацаа, жилийн хэлбэлзэл, Sharpe-ийн
харьцаа, багцын эргэлт, HHI үзүүлэлтийг
ижил аргачлалаар тооцно. Тэнцүү жинтэй
багцын төвлөрөл нь бүх сард
$HHI^{EW}=1/N$ байна.

Арилжааны зардлыг тооцохдоо нэг талын
арилжааны зардлын хувь хэмжээг $c$ гэж
тэмдэглэж, $c\in\{0,5,10,25\}$ суурь
пункт (basis points) гэсэн хувилбаруудыг
авч үзнэ. Энд нэг суурь пункт нь
$0.0001$-тэй тэнцэнэ.

Өмнөх сарын өгөөж хэрэгжсэний дараах
жин $\widetilde{\boldsymbol{w}}_t^{(k)}$
болон шинэ зорилтот жингийн ялгаагаар
тооцсон $TO_t^{(k)}$ нь нийт худалдан
авалт эсвэл нийт борлуулалтын эзлэх хувийг
илэрхийлнэ. Тиймээс нэг талын зардлын
хувь $c$ үед хоёр талын нийт арилжааны
зардлын хувь нь
\[
C_t^{(k)}=2c\,TO_t^{(k)}
\]
байна. Тухайн сарын эхэнд зардлыг
суутгасан гэж үзвэл цэвэр өгөөж нь
\[
r_{p,t}^{(k),net}
=\left(1-C_t^{(k)}\right)
\left(1+r_{p,t}^{(k)}\right)-1,
\]
энд $k\in\{0,1,\ldots,5,EW\}$.

OOS эхний сарын багцыг анх бүрдүүлэх
арилжааны зардлыг тооцохгүй, өөрөөр
хэлбэл $C_{T+1}^{(k)}=0$ гэж авна.
Үлдсэн 51 сарын дахин тэнцвэржүүлэлтийн
зардлыг тооцно. Арилжааны зардлын
хувилбар бүрд цэвэр өгөөжийн цуваагаар
жилийн хэлбэлзэл, Sharpe-ийн харьцаа,
хуримтлагдсан өгөөжийг дахин үнэлнэ.
Хуримтлагдсан цэвэр өгөөжийг
\[
CR_k^{net}
=\prod_{t=T+1}^{n}
\left(1+r_{p,t}^{(k),net}\right)-1
\]
гэж тодорхойлно.

% ----------------------------------------------------------
% 2.6.9. Bootstrap Comparison with the 1/N Benchmark
% ----------------------------------------------------------

\subsubsection{GMVP ба тэнцүү жинтэй багцын
	bootstrap харьцуулалт}
\label{subsubsec:bootstrap_equal_weight}

Үндсэн High, Medium, Low багц тус бүрд
зургаан GMVP загварын OOS вариацыг
$1/N$ багцын вариацтай харьцуулна.
$k$-р GMVP загварын $1/N$ багцтай
харьцуулсан дисперсийн бууралтыг
\[
RVR_{k\mid EW}^{(\%)}
=100\left(
1-\frac{\widehat{\sigma}_{p,k}^{\,2}}
{\widehat{\sigma}_{p,EW}^{\,2}}
\right),\qquad k=0,\ldots,5
\]
гэж тодорхойлно. Энэ үзүүлэлтийн
суурь нь өмнөх хэсгийн Sample GMVP
биш, харин $1/N$ багц юм.

Тодорхой бус байдлыг үнэлэхдээ
52 сарын OOS өгөөжид зургаан сарын
блокийн урттай circular moving-block
bootstrap аргыг хэрэглэнэ. Нэг багцын
долоон стратегийн (зургаан GMVP болон
$1/N$) өгөөжийг нэг ижил сарын
индексээр дахин түүвэрлэж, стратегиудын
хоорондын хос хамаарлыг хадгална.

Нийт $B=3000$ давталт бүрд түүврийн
вариац ($M-1$ хуваарьтай) болон
$RVR_{k\mid EW}^{(\%)}$-ийг шинээр
тооцно. Bootstrap тархалтын $2.5$ ба
$97.5$ хувийн квантилуудаар засвар
хийгдээгүй 95\%-ийн percentile итгэх
завсар байгуулна.

Мөн тэг таамаглал
$H_0:RVR_{k\mid EW}^{(\%)}=0$-ийг
хоёр талт, ажиглагдсан статистикийн
эргэн тойронд төвлөрүүлсэн bootstrap
ойролцооллоор үнэлнэ. Ажиглагдсан
RVR-ийг $\widehat{\theta}$, bootstrap
давталтын RVR-ийг $\widehat{\theta}^{*(b)}$
гэвэл тооцоолсон $p$-утга нь
\[
\widehat{p}_{boot}
=\frac{1+\sum_{b=1}^{B}
	\boldsymbol{1}\!\left\{
	\left|\widehat{\theta}^{*(b)}-\widehat{\theta}\right|
	\geq\left|\widehat{\theta}\right|
	\right\}}{B+1}.
\]

Нэг үндсэн багц дахь зургаан GMVP--$1/N$
харьцуулалтын $p$-утгуудад Holm-ийн
засвар хэрэглэнэ. Гурван багцын
нийт 18 шалгалтыг нэг бүл болгон
засварлахгүй. Percentile итгэх завсар
болон төвлөрүүлсэн bootstrap $p$-утга
нь нэг шалгалтын харилцан урвуу
хувилбар биш тул тэдгээрийн дүгнэлт
зарим тохиолдолд зөрж болно.

Энэ шинжилгээ нь зөвхөн үндсэн гурван
багцын зардлын өмнөх OOS өгөөжид
тулгуурлана. Иймээс арилжааны зардлын
дараах гүйцэтгэлийн статистик шалгалт
гэж тайлбарлахгүй. OOS хугацаа 52 сар
болохыг үр дүнгийн тодорхой бус байдлыг
тайлбарлахдаа харгалзан үзнэ.